% ============================================================
% CITY OF BRUTES — A Screenplay
% Author: Flyxion
% Compiled with LuaLaTeX
% ============================================================

\documentclass[12pt, oneside]{article}

% --- Fonts and Encoding ---
\usepackage{fontspec}
\usepackage{microtype}

% --- Multilingual Support ---
\usepackage{polyglossia}
\setmainlanguage{english}
\setotherlanguage{arabic}
\setotherlanguage{latin}
\newfontfamily\arabicfont[
  Script=Arabic,
  Scale=1.1,
  Path=/usr/share/fonts/truetype/noto/,
  Extension=.ttf,
  UprightFont=NotoNaskhArabic-Regular,
  BoldFont=NotoNaskhArabic-Bold
]{NotoNaskhArabic}

% --- Page Geometry ---
% Industry-standard screenplay margins
\usepackage[
  top=1in,
  bottom=1in,
  left=1.5in,
  right=1in,
  paperwidth=8.5in,
  paperheight=11in
]{geometry}

% --- Spacing ---
\usepackage{setspace}
\setstretch{1.15}

% --- Mathematics ---
\usepackage{amsmath, amssymb, amsfonts}

% --- Color ---
\usepackage[dvipsnames]{xcolor}
\definecolor{scenecolor}{gray}{0.08}
\definecolor{rulecolor}{gray}{0.45}

% --- Section and Heading Formatting ---
\usepackage{titlesec}
\titleformat{\section}
  {\large\bfseries\scshape\centering}
  {}{0em}{}
  [\vspace{0.2em}{\color{rulecolor}\hrule height 0.4pt}\vspace{0.4em}]
\titleformat{\subsection}
  {\normalsize\bfseries}
  {}{0em}{}

% --- Screenplay-Specific Environments and Commands ---

% Scene heading: all-caps, bold, with rule beneath
\newcommand{\sceneheading}[1]{%
  \vspace{1.8em}%
  \noindent{\bfseries\large\MakeUppercase{#1}}%
  \vspace{0.6em}%
  \par\noindent{\color{rulecolor}\rule{\linewidth}{0.3pt}}%
  \vspace{0.5em}\par%
}

% Action block: full-width, extended prose register
\newenvironment{action}{%
  \vspace{0.6em}%
  \noindent\ignorespaces%
}{%
  \par\vspace{0.6em}%
}

% Character cue: centered, all-caps, bold
\newcommand{\character}[1]{%
  \vspace{0.8em}%
  \begin{center}%
    \bfseries\MakeUppercase{#1}%
  \end{center}%
  \vspace{-0.3em}%
}

% Parenthetical: indented, italic
\newcommand{\paren}[1]{%
  \begin{center}%
    \itshape (#1)%
  \end{center}%
  \vspace{-0.3em}%
}

% Dialogue block: indented both sides
\newenvironment{dialogue}{%
  \vspace{0.2em}%
  \begin{list}{}{%
    \setlength{\leftmargin}{1.2in}%
    \setlength{\rightmargin}{0.8in}%
    \setlength{\itemsep}{0.6em}%
    \setlength{\parsep}{0pt}%
  }%
}{%
  \end{list}%
  \vspace{0.4em}%
}

% Dialogue item shorthand
\newcommand{\says}[2]{%
  \character{#1}%
  \begin{dialogue}\item #2\end{dialogue}%
}

% Parenthetical dialogue
\newcommand{\saysp}[3]{%
  \character{#1}%
  \paren{#2}%
  \begin{dialogue}\item #3\end{dialogue}%
}

% State marker: formal section break between states
\newcommand{\statebreak}[2]{%
  \vspace{2.5em}%
  \begin{center}%
    {\large\bfseries\scshape State #1: #2}%
  \end{center}%
  \vspace{0.3em}%
  \begin{center}%
    {\color{rulecolor}\rule{0.6\linewidth}{0.5pt}}%
  \end{center}%
  \vspace{1.5em}%
}

% Scene number marker
\newcommand{\scenenumber}[1]{%
  \vspace{1em}%
  \noindent{\small\bfseries\scshape Scene #1}%
  \vspace{0.3em}\par%
}

% Spherepop annotation: set in smaller italic type as marginal gloss
\newcommand{\spnote}[1]{%
  \vspace{0.4em}%
  \noindent{\small\itshape\color{rulecolor} #1}%
  \vspace{0.4em}\par%
}

% Transition
\newcommand{\transition}[1]{%
  \vspace{0.8em}%
  \begin{flushright}%
    \bfseries #1%
  \end{flushright}%
  \vspace{0.5em}%
}

% --- Lists ---
\usepackage{enumitem}
\setlist{nosep}

% --- Hyperlinks ---
\usepackage{hyperref}
\hypersetup{
  colorlinks=true,
  linkcolor=black,
  urlcolor=black,
  pdfauthor={Flyxion},
  pdftitle={City of Brutes}
}

% --- Header/Footer ---
\usepackage{fancyhdr}
\pagestyle{fancy}
\fancyhf{}
\fancyhead[R]{\small\itshape City of Brutes}
\fancyhead[L]{\small\itshape Flyxion}
\fancyfoot[C]{\small\thepage}
\renewcommand{\headrulewidth}{0.3pt}

% ============================================================
\begin{document}
% ============================================================

\thispagestyle{empty}

\vspace*{3em}

\begin{center}
  {\Huge\bfseries CITY OF BRUTES}\\[1.2em]
  {\large\itshape A Screenplay}\\[2.5em]
  {\normalsize Flyxion}\\[1em]
  {\small\itshape Independent Researcher}\\[3em]
\end{center}

\begin{center}
  {\small\color{rulecolor}\rule{0.5\linewidth}{0.4pt}}
\end{center}

\vspace{1.5em}

\begin{center}
\begin{minipage}{0.78\linewidth}
\small\itshape
\noindent\textit{Quod animalia bruta ratione utantur melius homine} --- that the brute animals employ reason better than man --- constitutes not a provocation but a hypothesis pressed upon its author by decades of irresolvable deliberation, a residue crystallized at the terminus of an institutional career that accumulated impossibility as its primary product. This screenplay is the event-history of that accumulation: not the treatise itself, but the sequence of irreversible states whose only admissible remainder was the treatise.
\end{minipage}
\end{center}

\vspace{2em}

\begin{center}
  {\small\color{rulecolor}\rule{0.5\linewidth}{0.4pt}}
\end{center}

\newpage

% ============================================================
% PREFACE
% ============================================================

\thispagestyle{empty}

\vspace*{2em}

\begin{center}
  {\large\bfseries\scshape Preface}
\end{center}

\vspace{1.5em}

\begin{spacing}{1.15}

This document does not present a narrative in the conventional sense, nor does it advance a thesis susceptible to demonstration through ordered exposition. It is, rather, the residue of a procedure: a series of selections made within a field of concurrent admissibilities, none of which, taken individually, possesses the authority to exclude the others. What appears here as continuity is not the unfolding of a single line but the compression of many, whose prior divergence is no longer directly recoverable from the surface that now presents itself.

\medskip

The scenes have been arranged without the guarantee of mutual stabilization. Each maintains local coherence --- its gestures, utterances, and configurations remain intelligible within their immediate frame --- yet their relations, when considered across the whole, resist consolidation into a single governing account. The work proceeds under the condition that contradiction does not necessarily entail breakdown, and that systems may persist, and even function, in the absence of global agreement.

\medskip

Three domains recur throughout: the administrative, the animal, and the performative. These are not to be understood as discrete categories but as distinct modalities through which order is enacted and maintained. The administrative operates through inscription, reference, and the progressive layering of directives whose authority derives from their continued circulation. The animal proceeds through distributed adjustment, without recourse to representation or archive. The performative introduces reiteration with variation, in which roles, claims, and identities detach from fixed bearers and circulate within structured repetition. No domain resolves the others. Each offers a partial articulation of a condition that exceeds it.

\medskip

Certain narratives appear in attenuated or displaced form: an offering accepted or refused without consensus; an act completed, prevented, or both; a mark whose function cannot be stabilized across those who encounter it. These do not serve as allegories to be decoded, nor as doctrinal positions to be affirmed or rejected. They persist as configurations --- reproducible, adaptable, and resistant to final adjudication. Their recurrence across different registers does not produce agreement, but a distribution of possibilities that remain in play.

\medskip

The language of the document reflects the condition it describes. Phrases recur --- \textit{as previously established}, \textit{as recognized}, \textit{as amended} --- without securing the references they invoke. Terms migrate between contexts, preserving syntactic function while altering their relational import. The reader may encounter sequences in which the local clarity of each statement does not extend to their cumulative relation. This is not a defect to be corrected but a property of the system being presented.

\medskip

In its final movements, the work undergoes a contraction. The multiplicity of prior forms is not resolved but selected from. Alternatives are neither refuted nor retained in parallel; they are excluded without commentary. What remains is a continuous surface whose legibility is immediate, though its formation is not. The absence of visible divergence should not be taken as evidence of prior agreement, but as the effect of a process that no longer displays the conditions of its own construction.

\medskip

No instruction is provided for the correct reading of what follows. The document does not require that its elements be reconciled, only that they be attended to as they appear. If coherence is sought, it must be constructed by the reader from within the constraints of what is present, without appeal to materials that are no longer accessible.

\end{spacing}

\vspace{2em}

\begin{center}
  {\small\color{rulecolor}\rule{0.4\linewidth}{0.4pt}}
\end{center}

\newpage

\noindent{\itshape The system operates at maximum channel density. Every institutional procedure functions as designed, which is to say it produces local coherences that fail to extend beyond their immediate radius. The human verbal channel carries everything. The perceptual logic of the film is present at its lowest admissible amplitude.}

\vspace{1.5em}

% ------------------------------------------------------------
\scenenumber{1}
% ------------------------------------------------------------

\sceneheading{INT. Diplomatic Antechamber — Day}

\begin{action}
The room presents itself first as an enclosure of moderate dimensions, plastered and irregularly lit, the light entering from a single elevated window and diffusing across the interior with the mild incoherence characteristic of early morning illumination when it has not yet found its angle of incidence upon any surface long enough to constitute a stable relationship with it. The walls, which had perhaps at some earlier administrative moment borne maps, portraits, or the certificates of delegated authority, retain only the faint discolorations of these former presences --- rectangular lacunae in the plaster's patina, slightly paler than their surrounds, which together constitute a kind of negative cartography of the room's prior semiotic life. What remains is a surface that has been subtracted from, not one that has been composed.

Several groups occupy the room simultaneously, their arrangement governed not by any overarching principle of spatial organization but by the accumulated residue of prior movements --- arrivals, exchanges, partial withdrawals --- such that the configuration of bodies at this moment is less a deliberate disposition than the current state of a system that has been in operation for some time and will continue in operation after the present observation concludes. No single arrangement governs their placement. Each cluster occupies the spatial interval that was available when its members arrived, and has since been modified by the slow deformations that prolonged proximity introduces.
\end{action}

\begin{action}
Voices overlap, not loudly but continuously, in the manner of a room that has achieved a certain acoustic equilibrium --- one in which speech has become ambient rather than directed, a medium rather than a message. The voices constitute a distributed field of communicative activity whose individual threads are locally coherent but globally irresolvable into any unified semantic structure. One does not listen so much as inhabit the sound.

At the center table: papers, folded and refolded, bearing the marks of successive handling. Some carry intact seals, suggesting they have not yet been opened; others bear the residue of seals broken and, in certain cases, imperfectly reapplied, their wax spread with insufficient pressure across a seam that no longer closes with authority. Still others have been left open, their contents available to any eye that passes over them, which is to say their confidentiality has been dissolved not through any decisive breach but through the accumulated indifference of an environment in which private information circulates too rapidly to retain the conditions of its privacy.
\end{action}

\begin{action}
\textsc{Rorario} occupies a position slightly off-center within the room --- neither fully engaged in any one conversational cluster nor withdrawn to its margins. His situation within the space is best understood as a function of local pressure rather than deliberative positioning: he has arrived at his current location through a series of minor accommodations to the movements of others, each adjustment individually negligible, their cumulative vector having produced a displacement from any center he might have originally intended to occupy. He stands with the practiced stillness of a man accustomed to receiving information from multiple directions simultaneously, his orientation shifting through small arcs as different sources of speech and gesture momentarily acquire precedence.

His face exhibits the trained illegibility of career diplomatic practice --- not blankness, which would be its own species of expression, but something closer to the suspension of expressivity as a default condition, a surface that has learned to receive without reflecting. He is, in the vocabulary of early modern institutional life, a perfectly functional node in a network of information exchange whose overall topology he can model but not modify.
\end{action}

\begin{action}
At the far end of the room, two \textsc{Diplomats} maintain a conversation whose formal register is that of collegial exchange and whose actual content is something considerably more adversarial, the adversarialism expressed entirely through modality --- through the selection of verb tenses and the placement of qualifiers --- rather than through any direct contestation of fact.
\end{action}

\says{First Diplomat}{The courier departed before dawn. The road was clear.}

\saysp{Second Diplomat}{after a measured interval}{That is not what we were told.}

\says{First Diplomat}{The report you received was compiled yesterday.}

\says{Second Diplomat}{The report I received was subsequently revised.}

\begin{action}
Neither voice rises. The contradiction persists between them in a state of suspension that is neither acknowledged nor dissolved, subsisting as a kind of shared premise upon which they continue to erect the structure of their exchange, each statement implicitly granting the prior statement a provisional validity it has not actually established. The conversation proceeds through contradiction without addressing it --- a phenomenon less extraordinary, in the context of this room and this institution, than it might elsewhere appear.
\end{action}

\begin{action}
Adjacent to the central table, a \textsc{Scribe} attends to the arrangement of documents with a concentration that has the outward form of purposefulness but whose actual trajectory is one of iterative micro-correction without convergence. He aligns the edges of a stack with care, examines the result, finds in it some insufficiency not immediately apparent to an external observer, and realigns. A seal on one document cracks under the slight pressure of his handling --- not broken, but compromised, its integrity now conditional upon the gentleness of all future contact with it. He examines the damage for a moment with the professional equanimity of someone who has learned that such events are neither rare nor consequential within their proper context, then places the document atop a different stack, reclassifying it according to a taxonomy that the external viewer cannot reconstruct from available evidence.
\end{action}

\begin{action}
A \textsc{Servant} crosses the room bearing a shallow tray upon which three cups have been arranged. The liquid within them trembles with the fine vibration of ambulation --- a surface-tension phenomenon that registers the movement of the carrier in its own medium and communicates it upward as a kind of somatic transcription of transit. He moves toward \textsc{Rorario} with the tentative trajectory of someone navigating a spatial field whose social geometry is not fully legible to him, then, arriving within a meter of his target, redirects --- placing the tray instead upon the table near the two Diplomats, who do not acknowledge the delivery, which is to say they receive it as a function of the room's service infrastructure rather than as a gesture requiring response.

\textsc{Rorario}'s hand lifts slightly --- the initiatory phase of an intercepting gesture --- and then descends, the intention remaining unexecuted, the motion folded back into stillness before it could declare itself.
\end{action}

\begin{action}
At the room's periphery, partially occluded by the lateral projection of a heavy bench, a \textsc{Dog} reclines upon the flagstone floor. Its attention, however, is not peripheral: its eyes conduct a continuous low-amplitude survey of the room's activity with the unhurried precision of a sensory apparatus calibrated to detect changes in the distribution of mass and motion rather than in the content of speech. When the Servant places the tray, the Dog rises without urgency, performs a single rotational movement of economical radius, and resettles itself with its dorsal axis now oriented toward the room's center. The reorientation is not a response to any legible stimulus; it is an adjustment, and adjustments in the Dog's behavioral repertoire have a quality of inexorability that distinguishes them sharply from the hesitative, self-revising movements characteristic of the room's human occupants.
\end{action}

\begin{action}
A \textsc{Translator} positions himself between two conversational clusters, functioning as a relay through which utterances generated in one linguistic register are converted into approximate equivalents in another. The conversion is performed with minimal delay and moderate fidelity --- the structural content of each statement is preserved while its modal texture undergoes systematic simplification, the qualifications and hedges of the original flattened into declarative forms whose apparent confidence exceeds that of the source material.
\end{action}

\says{Voice (O.S.)}{The terms were accepted in full.}

\says{Translator}{The terms were accepted.}

\says{Voice (O.S.)}{In full.}

\saysp{Translator}{after a hesitation of roughly two seconds}{The terms were accepted without registered objection.}

\begin{action}
The substitution of \textit{without objection} for \textit{in full} effects a semantic displacement of some significance: the first formulation designates a positive condition, the second merely the absence of a negative one. Whether this distinction is appreciable to the recipients of the translation is not disclosed by their response, which is to receive the formulation and proceed, their own subsequent statements now resting upon a premise that has been slightly but irreversibly altered from the one upon which the original speaker was operating.

\textsc{Rorario} moves incrementally toward the Translator's position. He does not attend exclusively to the translation; his attention distributes itself across the room in a pattern of continuous partial engagement, touching each source of information briefly before relocating. His gaze traverses speakers, documents, the arrangement of the tray, the Dog's new orientation.
\end{action}

\begin{action}
A second exchange initiates itself in the room's posterior region, partially overlapping with the first in both time and acoustic space.
\end{action}

\says{Voice (O.S.)}{The envoy has already arrived.}

\says{Second Voice (O.S.)}{That is not possible.}

\says{Voice (O.S.)}{He was observed at the eastern gate.}

\says{Second Voice (O.S.)}{Then the figure at the eastern gate was not him.}

\begin{action}
The Translator, perceiving this new thread at the periphery of his operational radius, executes a slight rotation of the shoulders --- the initiatory gesture of reorientation --- but does not abandon the relay he is currently conducting. The two conversations proceed in parallel, their intersection constituted not by any shared informational content but by the shared acoustic field they occupy and the single body of the Translator who is insufficiently multiplied to serve them simultaneously without loss.

\textsc{Rorario} reaches toward the central table and selects a document --- the selection has the form of deliberation without its substance, the hand arriving at a particular document through a movement that is neither random nor clearly motivated. He scans the uppermost portion of the text. His eyes pause upon one line --- the pause briefer than a reading pause, longer than a glance --- and then continue to the document's lower portion without extracting the intermediate content. He returns the document to the table in a position displaced some four centimeters from where it lay, a displacement insufficient to constitute a reorganization and sufficient to introduce a small local perturbation into the Scribe's ongoing rectification project.

The Scribe perceives the displacement, adjusts the stack, pauses in a condition of minor underdetermination regarding which configuration more closely approximates the arrangement he was attempting to restore, and leaves the matter as it stands.
\end{action}

\begin{action}
The Servant returns, his tray now emptied of its cups, and navigates again toward \textsc{Rorario} --- this time with the slight increase in directedness that second approaches tend to acquire when the first has been inconclusive. \textsc{Rorario} receives him.
\end{action}

\says{Rorario}{Was this delivered?}

\saysp{Servant}{with the careful imprecision of someone uncertain whether the question has a correct answer}{It was placed.}

\begin{action}
The distinction between \textit{delivered} and \textit{placed} is not pursued. \textsc{Rorario} acknowledges the response with a slight inclination of the head --- the gesture of receipt rather than assent, accepting the statement into the communicative record without conferring upon it the status of a satisfactory answer to the question that occasioned it. The gap between question and response remains, folded into the exchange's apparent completion.

Across the room, one of the Diplomats lifts a cup from the tray and sets it down again without having brought it to his lips, the gesture describing a complete arc of intention without executing the purpose that arc implied. The Dog, which has been stationary since its earlier reorientation, now rises again, crosses the room on a trajectory that brings it to within a half-meter of the central table, and arrests its forward motion at this point --- not, it appears, due to any obstacle or inhibitory stimulus, but simply because the trajectory has arrived at a natural terminus that the animal accepts without requiring it to correspond to any externally visible destination. It looks upward briefly --- the angle of the gaze suggesting attention to something at approximately human chest height, though no particular feature of the room at that elevation presents itself as an obvious attractor --- and then executes a turn and exits through the partially open door with the characteristic unhesitating continuity that marks each of its movements: not fast, but without the interpolated pauses of revision that interrupt the locomotion of everyone else in the room.
\end{action}

\begin{action}
\textsc{Rorario} observes the door as it comes to rest in its partially closed position. The exterior light that admits itself through the gap is of a quality and intensity sufficient to indicate the existence of an outside but insufficient to disclose its particular character. The door does not fully close; it does not fully open. It maintains its position in a condition of unresolved mediation between two states.

Behind him, the Translator's voice continues its relay function, now at a slight reduction in semantic density relative to its earlier iterations.
\end{action}

\says{Translator}{The terms remain unchanged.}

\begin{action}
The formulation, which has passed through several prior iterations each of which introduced some small modification of emphasis or qualification, has arrived at this form --- its most reduced and apparently stable --- not through a process of deliberate distillation but through the accumulated effects of successive transmission. It is not a summary; it is what remains after summary has been applied repeatedly to material that each iteration rendered slightly more attenuated. That it now sounds authoritative is an artifact of its compression, not an index of its fidelity to whatever it originally designated.

\textsc{Rorario} gathers two documents from the table --- one bearing an intact seal, one already opened and refolded along imprecise creases. He does not examine them in relation to each other. He folds the opened document once more, pressing along its existing creases to deepen them without improving their alignment, and inserts it within the folded exterior of the sealed one, pressing the composite at its edges without fastening it --- an act of combination that produces the superficial appearance of a unified object while leaving its internal relation entirely undecided, the outer sheet's formal closure enclosing an inner sheet's prior breach without annulling it.

He holds this composite object between his palms for a moment, applying slight lateral pressure as if testing its coherence. The composite offers no resistance and provides no confirmation. \textsc{Rorario} releases it, returns it to the table, and withdraws from the table's immediate vicinity.

The conversations continue in the room's ambient register. The Scribe continues his project of micro-rectification. The Servant has departed the scene. The door through which the Dog exited stands at its indeterminate angle.
\end{action}

\begin{action}
What constitutes the phenomenological substance of this scene is neither the content of any particular exchange nor the resolution of any of the several contradictions that have been introduced and left suspended within it, but rather the specific quality of cognitive labor that the room imposes upon anyone who attempts to occupy it as a unified epistemic subject --- the continuous and ultimately unsuccessful effort to maintain commensurability between competing streams of information that do not share a common frame of reference, do not arrive at a common rate, and are not addressed to any single receiver whose synthetic capacity might be equal to their simultaneous demands. \textsc{Rorario} does not fail in this effort; he is, by every available measure, succeeding at the highest level of institutional competence. What he produces, however, is not comprehension but its functional substitute: a continuous stream of locally adequate responses whose global incoherence is not registered as such because the system within which they circulate has no mechanism for doing so.

The composite document on the table, its sealed exterior enclosing its already-opened interior, is this situation's material emblem --- though it is not presented as such, since the film does not present emblems. It is simply an object that has been produced by the room's operations and will continue to exist within it, its internal irresolution persisting beneath a surface that appears, to any external examination, to have been concluded.
\end{action}

\transition{CUT TO:}

% ------------------------------------------------------------
\scenenumber{2}
% ------------------------------------------------------------

\sceneheading{INT. Private Reception Chamber — Day}

\begin{action}
The corridor has deposited \textsc{Rorario} into a space of different acoustic and luminous character --- smaller, its proportions more deliberate, its illumination originating from two sources whose angles of incidence produce a mild interference pattern of overlapping shadows across the central floor. The room has the quality of having been prepared, which is to say it has been recently arranged by someone whose purpose was to create the impression of natural order rather than to produce actual functional organization; the result is a space in which everything is slightly more symmetrical than use would have made it and slightly less symmetrical than overt deliberation would have achieved.

A \textsc{Noblewoman} occupies the room as its primary figure. She is seated in a chair positioned at a slight angle to the room's axis of symmetry --- an angle produced either by her habitual preference or by the room's imperfect geometry, and one which gives her the aspect of someone who has settled into the space rather than been placed within it. Her posture is neither formally erect in the manner of ceremonial reception nor casually disposed in the manner of private repose; it is the posture of a person for whom the distinction between these modes is not currently operative, because she is engaged in the functional work of the meeting rather than its performative dimensions.
\end{action}

\begin{action}
\textsc{Rorario} enters and is received with the minimal protocol appropriate to a meeting whose official status is ambiguous --- not formal enough for the full ceremonial register, not informal enough to dispense with ceremony altogether. The exchange of greetings is handled with the economy of two people who understand their relative positions without requiring them to be articulated.

The meeting proceeds. Its ostensible subject is the disposition of certain territorial questions relating to ecclesiastical jurisdiction over a district whose administrative status has become contested through a concatenation of events each of which was locally resolvable but whose combination has produced a situation that the existing frameworks of both secular and religious law address imperfectly and in partially incompatible ways. \textsc{Rorario} has come with a position; the Noblewoman has a position. The positions are not opposed in any simple sense; they are differently shaped, proceeding from different premises, addressing different aspects of the same underlying configuration, and capable of being brought into approximate alignment through the kind of careful mutual adjustment that constitutes the practical substance of negotiation at this level.
\end{action}

\begin{action}
The Noblewoman conducts the meeting with a quality of attention that is, in certain respects, more sophisticated than the institutional register typically demands. She does not merely receive \textsc{Rorario}'s formulations and respond to them; she tracks the relationship between successive formulations, noting where they are consistent and where they introduce small modifications that may or may not be intentional. She does not call attention to these observations; she incorporates them into the structure of her own responses, which are consequently more precisely targeted than they would be if she were responding only to the surface content of each statement rather than to its trajectory.

This is not described because it is unusual --- it is, in fact, a common capacity among those who practice negotiation with sustained attention --- but because it is not immediately visible and will not be registered by anyone in the room as worthy of remark. The Noblewoman's analytical precision subsists as an unacknowledged operative force within the encounter, shaping its outcomes without entering its explicit record.
\end{action}

\begin{action}
An agreement of the approximate kind that these meetings are designed to produce begins to assume legible form. Its terms are not ideal from either party's perspective --- ideal outcomes are not available in the constraint environment within which both parties operate --- but they are admissible, which is the operative criterion. \textsc{Rorario} formulates the emerging consensus in a condensed Latin phrase whose compression performs the work of making the approximate appear precise. The Noblewoman receives this formulation, considers it through an interval of roughly four seconds in which her gaze does not leave his, and offers a modification of one term --- a modification small enough to leave the overall structure intact and significant enough to constitute a genuine adjustment of emphasis.

\textsc{Rorario} accepts the modification. The agreement has assumed its form.
\end{action}

\begin{action}
The \textsc{Husband} enters.

He enters from a lateral door whose existence had not previously been indicated by the room's visible architecture --- a door that presents itself, in the moment of opening, as both an architectural fact and a narrative one, its opening coinciding with the agreement's completion in a manner that cannot be described as fortuitous without implying a degree of design that the scene does not permit.

He is a man of middle years whose physical presence in the room immediately alters its social geometry in ways that are evident without being easily specifiable. He is not loud, not imposing, not hostile. He is simply a presence whose authority within this domestic-institutional space predates and supersedes the provisional authority generated by the meeting that has just concluded, and this supersession is felt immediately by both parties without being declared.
\end{action}

\says{Husband}{The question of the northern boundary has already been referred to the Cardinal's office. Any agreement reached here would be subject to that prior referral.}

\begin{action}
The statement does not contest the agreement directly; it contextualizes it in a way that, if accurate, renders it provisional rather than final --- a distinction that is, in practice, equivalent to dissolution. Whether the referral to the Cardinal's office is accurate, recent, or operationally relevant is not verifiable within the room. The Husband delivers the statement with the mild certainty of someone reporting an administrative fact rather than advancing a negotiating position.

He nods --- a gesture of acknowledgment addressed to \textsc{Rorario} that simultaneously functions as a conclusion --- and exits through the same lateral door by which he entered, which closes behind him with a completeness that the antechamber's door did not achieve.

The room reassumes its prior configuration.
\end{action}

\begin{action}
\textsc{Rorario} and the \textsc{Noblewoman} regard each other across the space of the just-concluded meeting.

The Noblewoman speaks first.
\end{action}

\says{Noblewoman}{I will send word when the referral has been clarified.}

\begin{action}
The formulation is perfect in its operational precision: it neither acknowledges the agreement's effective dissolution nor pretends the meeting has produced a stable outcome; it designates a future communicative act that will occur when conditions permit and that will, at that point, resume the meeting's suspended work. It treats the interruption as an administrative interval rather than a substantive reversal.

\textsc{Rorario} receives this with a slight inclination of the head identical in form to the gesture he directed at the Servant in the prior scene --- the gesture of receipt without conferral of assent, acknowledging the statement's entry into the record without endorsing its implicit characterization of what has occurred.

He gathers his documents. He has not written anything during the meeting. What had been an emerging agreement exists now only in the memories of the two parties and in the ambient record of the room, which retains nothing. He takes his leave with the formal brevity appropriate to a meeting whose conclusion is ambiguous and whose parties have agreed, without discussion, to treat the ambiguity as a form of provisional closure.

The room, as he exits, returns to the condition of prepared stillness in which he found it, the Noblewoman remaining in her chair at its slight angle to the room's axis, the lateral door closed, the agreement's residue distributed with perfect evenness between presence and absence --- present enough to have occurred, absent enough to require repetition.
\end{action}

\transition{CUT TO:}

% ------------------------------------------------------------
\scenenumber{3}
% ------------------------------------------------------------

\sceneheading{INT. Formal Dining Hall — Evening}

\begin{action}
The hall presents itself under the aspect of designed abundance: a table of considerable length set with the careful asymmetry that distinguishes genuine use from ceremonial arrangement, its surface bearing the accumulated evidence of a meal that has been in progress for some time --- platters partially depleted, cups at various stages of consumption, the residue of earlier courses not yet fully cleared, which produces a stratigraphy of the meal's own history visible in the overlapping rings left by vessels, in the arrangement of displaced objects, in the particular geometry of spaces opened by consumption within the table's original arrangement. The light comes from candles distributed at irregular intervals, each casting its own local illumination whose overlap with neighboring candles produces zones of neither full brightness nor full shadow but of a diffuse, softly contested luminosity that renders faces with the slightly indeterminate quality of objects seen in conditions that do not fully commit to either visibility or obscurity.

The room is occupied by perhaps fourteen figures, their seating arrangement encoding a social hierarchy that has already been partially disordered by the progressive relaxation of formality that a meal of some duration necessarily introduces. Voices overlap in the manner now familiar from the antechamber, but with the added modality of food and drink --- a different quality of simultaneous utterance, less transactional and more ambient, though not less continuous.
\end{action}

\begin{action}
\textsc{Rorario} is seated at a position of moderate eminence within the table's hierarchy --- not at its head, which would imply a primacy he does not hold in this particular configuration, but sufficiently toward the center that the conversational field within his reach is broad and multiply sourced. He eats with the systematic inattention of someone for whom the meal is primarily an occasion for other activities, his fork and knife operating with a low-amplitude automatism that leaves the preponderance of his cognitive resources available for the room's communicative field.

To his left, a \textsc{Senior Diplomat} holds a position of slightly greater eminence, currently engaged in a detailed account of a negotiation concluded several months prior, an account whose relationship to the negotiation's actual proceedings is, as these retrospective narratives invariably are, one of selective emphasis and structural simplification --- the messy simultaneity of actual events having been converted, through the operations of memory and narration, into a clean sequential causality whose clarity is an artifact of retrospection rather than a feature of what occurred. His interlocutors receive the account with the attentive expressions of those who understand that the purpose of such narratives is not informational but social, and respond accordingly.
\end{action}

\begin{action}
The \textsc{Senior Diplomat} has been speaking for some time when he pauses mid-sentence --- the pause having the character of an intake of breath before continuation rather than of any terminal hesitation --- and then does not continue. He remains for a moment in a condition of arrested mid-utterance, his expression sustaining the communicative intention that preceded it but no longer producing its articulation, and then, with the abrupt economy of a system that has ceased operation without transition, he inclines forward and to the left, his shoulder making contact with the back of his chair before his body settles into an attitude of collapse that is less dramatic than such events are commonly represented and more simply administrative in its character --- a body that has stopped functioning in its current mode and requires redistribution within the space it occupies.

The \textsc{First Guest} to his immediate left receives this movement against his own shoulder before perceiving its cause, at which point he turns and assesses the situation with the practiced rapidity of someone who has encountered, if not this precise event, its functional category before. He does not exclaim. He makes a brief and quiet statement to the person on his other side, who relays it to the next, the information propagating down the table's length by a process of successive whispered transmission that is structurally indistinguishable from the communicative channels already operative in the room.

A \textsc{Servant} at the room's margin, observing the situation from the wider vantage of his peripheral position, begins to move toward the table but arrests this movement at approximately the midpoint of the room, having arrived at an underdetermination regarding the appropriate sequential relationship between the removal of the affected person and the removal of the current course, neither of which admits an obvious temporal priority over the other, and both of which fall within the range of his responsibilities. He remains in this position of suspended decision for the duration of approximately four seconds, then continues forward, his choice having been made through the accumulated pressure of the delay itself rather than through the resolution of the underlying underdetermination.
\end{action}

\begin{action}
\textsc{Rorario} has observed the Senior Diplomat's collapse with the lateral attention of someone who processes peripheral events without redirecting his primary focus. By the time the relay of whispered information has traveled far enough down the table to constitute a social acknowledgment of the event, \textsc{Rorario} has already, without announcement or visible transition, extended the conversational range of his own engagement to incorporate the interlocutors who had been addressed to the Senior Diplomat --- absorbing their attention into his own field of engagement with the quiet efficiency of a system reconfiguring itself in response to a change in available channels. He does not assume the Senior Diplomat's position; he does not move his seat or alter his physical location. He simply becomes, through a redistribution of his communicative output, the operative node in a network that has lost one of its elements and reorganized around the loss.

The conversations in his immediate vicinity continue without interruption. Several of those further from the event have not yet received the relay. The \textsc{Servant} at the room's margin has made his decision and acts upon it, the body of the Senior Diplomat being assisted toward a lateral exit with a discretion so practiced as to render the process nearly invisible within the room's ambient activity.

At the hall's far margin, partially occluded by the leg of a secondary serving table, a \textsc{Dog} lies in the attitude of repose. It has been in this location since before the scene's initiation, or at any rate since before the current viewer's attention was directed toward this particular corner of the room. Its position is, however, subtly different from the position it occupied at the scene's beginning --- an observation that becomes available only through retrospective comparison, through the recognition that the dog's dorsal axis, which had been oriented along one diagonal of the available space, is now oriented along a slightly different one, and that the distance between the dog's present location and the serving table's leg, which had been approximately twenty centimeters, is now approximately forty. The repositioning has occurred without being witnessed, without producing any audible indication of movement, without attracting the attention of any human occupant of the room. It is simply the case that the dog is now located where it is, and was previously located somewhere slightly different.

One \textsc{Guest}, in the middle of a statement whose subordinate clauses have accumulated to a degree that makes its syntactic conclusion genuinely uncertain, redirects his utterance mid-course in response to some stimulus not visible to the external viewer, producing a sentence that ends in a different grammatical mode than it began, with a conclusion that addresses a subject the opening had not implied. He does not mark the transition. His interlocutor receives the resulting statement without registering the disjunction as remarkable, incorporating it into the exchange's running record as if it were the utterance that had been intended from the beginning.

Someone, in the vicinity of the displaced cup now occupying the space vacated by the Senior Diplomat's tray, spills wine --- the spill modest in volume, its propagation across the tablecloth arrested by an adjacent fold of the linen before it can reach the documents that have accumulated near \textsc{Rorario}'s position. A second Servant moves to address the spill with the focused competence that minor physical mishaps in formal contexts elicit, applying cloth to the surface with the practiced pressure of someone who has determined through experience the precise degree of absorption required and the precise gesture that achieves it.

The table's activity continues in its ambient register.
\end{action}

\transition{CUT TO:}

% ------------------------------------------------------------
\scenenumber{4}
% ------------------------------------------------------------

\sceneheading{INT. Administrative Writing Room — Day}

\begin{action}
The room is functional in the undecorated sense --- a space whose architecture exists entirely in service of a specific category of activity and which has, over time, been stripped of any element not in direct relation to that activity, not through deliberate austerity but through the accumulated priority of use over representation. The walls are plain. The single window is positioned for light rather than for view. The furniture --- two chairs, a writing table, a secondary surface for the arrangement of documents and instruments --- has been placed according to the logic of workflow rather than of social geometry. The room is not austere in any spiritually intentional sense; it is simply a room in which the overhead of appearance has been reduced to approximately zero.

Two documents lie on the writing table, having been placed there by someone before the scene's commencement. They are sealed --- both of them --- with the impressions of different sigillary instruments, the distinctions between the two seals visible to anyone with sufficient familiarity with the relevant sigillary conventions to read them, which within this institution means: nearly everyone, and certainly \textsc{Rorario}.

He enters from the corridor, accompanied by a \textsc{Clerk} whose function is administrative rather than diplomatic --- a man whose competence operates within a narrower but more precisely defined range than \textsc{Rorario}'s own, and whose relationship to the documents on the table is one of custody rather than interpretation.
\end{action}

\begin{action}
The Clerk delivers a verbal summary of the documents' contents without opening them, demonstrating thereby a memorization of their substance that implies either prior access to their contents or a reconstruction from contextual knowledge sufficient to render prior access unnecessary. The summary is delivered with the practiced neutrality of an administrative register --- neither emphasizing nor minimizing, applying the same tonal weight to each element regardless of its apparent significance, a procedure that distributes the receiver's attention evenly across the whole rather than directing it toward any particular component.

The summary has been underway for perhaps ninety seconds when \textsc{Rorario} opens the first document.
\end{action}

\begin{action}
The document's contents, which he reads while the Clerk continues his verbal summary of the same document, occupy a register that is not identical to the summary's. Not contradictory --- the factual substance of the two accounts is substantially congruent --- but differently emphasized, differently structured, proceeding from different premises regarding what requires explicit articulation and what may be assumed. \textsc{Rorario} reads and listens simultaneously, the two streams of information arriving through different sensory channels at slightly different rates, producing a stereoptic epistemic effect in which the document and its summary cast slightly different shadows of the same object.

He turns the page. The Clerk continues.
\end{action}

\begin{action}
The second document, which \textsc{Rorario} opens before the Clerk's summary of the first has concluded, bears instructions from a source whose authority over the matter at hand exists in a relationship of uncertain precedence with the authority of the first document's source. The two documents have been issued by offices whose jurisdictional relationship is defined by a series of agreements and customary arrangements that were themselves produced through a process of accumulated negotiation over several decades, and whose coherence in any specific case depends upon an interpretive consensus that is itself subject to ongoing revision. In this case, the instructions they respectively issue address the same situation and arrive at positions that are, if not directly contradictory, then oriented along different axes --- each internally consistent, each locally authoritative, together constituting a configuration that has no single admissible resolution within the frameworks either document recognizes.

The Clerk, who has been aware of this situation since the documents arrived, delivers his summary of the second document with the same tonal neutrality he applied to the first.
\end{action}

\says{Clerk}{The instructions in the second document concern the same matter as the first. The positions differ on the question of temporal priority. Both offices have requested a response.}

\begin{action}
\textsc{Rorario} places both documents on the table, their seals now broken, their contents now inhabiting the same physical space though they continue to inhabit different interpretive ones. He regards them together for a moment --- the posture of someone who is not deliberating in any exploratory sense but is allowing the configuration to stabilize within his cognitive field before acting upon it, the way a painter regards a wet surface before committing to the next mark.

Then he takes up a pen, draws a sheet of paper, and begins to write.

He writes without visible hesitation, without pausing to re-read what he has written, without the iterative checking movements --- backward glance to source document, forward glance to developing text --- that characterize a writer uncertain of the relationship between his source material and his output. The text he produces is not a synthesis of the two documents in any dialectical sense; it is not a resolution of their differences, nor a subordination of one to the other, nor an acknowledgment of the tension between them addressed to some higher authority capable of adjudicating it. It is, rather, a formulation whose internal logic is self-consistent and whose relationship to both source documents is one of selective extraction --- taking from each the elements that do not conflict with the elements taken from the other, and constructing from this selection a statement that is coherent within itself and whose coherence does not require the incoherence of its sources to be addressed. The incoherence is absorbed into the production process and does not appear in the product.

The resulting document will be received by its addressees as a third position --- one that appears to proceed from the two documents while actually proceeding from the space between them, a space that the documents' authors did not design but inadvertently produced through the combined action of their individually coherent but mutually discordant instructions.
\end{action}

\begin{action}
\textsc{Rorario} sands the page, folds it, and seals it with his own instrument --- an act that transforms the document from a draft into a dispatch, from a private production into an institutional communication, from something that could still be revised into something that has been committed. The seal's wax cools. The fold holds. The document has acquired the formal properties of a completed act.

He hands it to the Clerk without annotation or instruction, the document's destination being implicit in the context of its production and requiring no supplementary specification.

The Clerk receives it, examines the seal with brief professional attention, and departs.

\textsc{Rorario} remains at the writing table. The two source documents lie open before him, their positions unchanged. He does not return them to their sealed condition; he does not stack them or otherwise reorganize them. He leaves them as they are --- open, their contents available, their contradictions intact --- and exits the room.

The documents remain on the table. The window continues to admit light. The room maintains its functional condition of preparation for the next use to which it will be put, indifferent to the particular use just concluded.

This is the terminal condition of State~I: the system at the apex of its operational fluency, producing coherent outputs from incoherent inputs through the application of an absorptive intelligence so practiced that the absorption is itself invisible, leaving no residue in the product that would indicate the difficulty of its production. Nothing has been eliminated yet in any visible sense; the options that have been narrowed by each of \textsc{Rorario}'s decisions appear, from the perspective of the decisions themselves, to remain open --- his documents go forth into a network that continues to generate new inputs at a rate sufficient to maintain the appearance of an open system. Only retrospectively will it become apparent that certain trajectories have already been foreclosed, their foreclosure having been performed through the very competence with which \textsc{Rorario} continues to function within the field that is, at this moment, still producing him as its most accomplished instrument.
\end{action}

\transition{CUT TO:}

% ============================================================
\statebreak{II}{Accumulation}
% ============================================================

\noindent{\itshape Three convergent pressure tracks initiate their parallel operations: the geopolitical non-termination of post-Moh\'acs Central Europe, the perceptual counterexample of animal action within human institutional space, and the structural incoherence of the papal administrative apparatus. These tracks do not yet compound; they coexist. The perceptual reweighting of signal --- the gradual redistribution of what the film treats as information --- begins its increment. Nothing breaks. Small incompatibilities begin to persist instead of dissipate.}

\vspace{1.5em}

% ------------------------------------------------------------
\scenenumber{3a}
% ------------------------------------------------------------

\sceneheading{EXT. Stable Courtyard — Morning}

\begin{action}
The stable courtyard occupies a position in the institutional complex that is adjacent to but formally distinct from the spaces of administration and reception --- a zone in which the logistical substrate of diplomatic life maintains itself with a visibility that the interior spaces work, through their arrangement and their protocols, to suppress. Here the relationship between deliberative human activity and the biological and mechanical infrastructure that sustains it is not mediated by the conventions of institutional representation; the two coexist in the unidealized proximity of functional necessity.

The yard is stone-floored, its surface worn into a topography of shallow depressions and low ridges by decades of hooved contact, the wear patterns recording the recurrent trajectories of animals moving between stable and open air, the historical accumulation of those movements legible in the stone's surface the way that paths become legible in ground that has been crossed many times. It is morning; the light enters the yard at an angle that throws the stone's surface variations into mild relief, each small depression gathering a slightly deeper shadow than its surroundings, such that the floor's topography is more visible now than it would be at midday, when direct overhead illumination would flatten the surface back into apparent uniformity.

A \textsc{Farrier} works at the yard's near margin, positioned beside a \textsc{Horse} that stands with the relaxed weight-distribution of an animal habituated to this procedure. The Farrier has the horse's right foreleg raised and held between his knees, the hoof presented upward for access. He works at its surface with a rasp, the motion of the tool repetitive and precisely modulated --- not identical strokes but variations within a narrow range, each adjusted by some feedback mechanism that the external observer cannot fully reconstruct, the adjustment made at the level of applied pressure and angle rather than at the level of deliberated decision.
\end{action}

\begin{action}
The \textsc{Horse}'s adjustments to the constraint of having one leg raised and held are continuous and of small magnitude --- a redistribution of weight across the three remaining points of contact with the ground, each shift compensating for some perturbation introduced by the Farrier's work or by the animal's own muscular activity in maintaining the raised posture. These adjustments do not interrupt the Farrier's work; they are too fine-grained and too smoothly executed to constitute events in any sense that would require response. They are, rather, a form of ongoing negotiation between the animal's proprioceptive field and the constraint imposed upon it, conducted entirely below the threshold of the deliberative, entirely within the domain of the immediately responsive.

The horse does not consider how to stand. It stands, continuously, making the thousand small corrections that standing on three legs while a fourth is worked upon requires, each correction performed at the speed the situation demands and with the precision the situation requires and without any of the interpolated pauses of self-assessment that would, in a human performing a comparable task, indicate the presence of executive oversight.

There is no hesitation in it. Not the suppression of hesitation, which would itself be a form of effort --- simply its absence, as a structural condition of this kind of action.
\end{action}

\begin{action}
\textsc{Rorario} passes through the stable yard in transit between two points whose institutional significance exceeds that of the yard itself. He is carrying a document. He slows as he moves through the Farrier's working radius, not stopping, but allowing the pace of transit to relax into something closer to an extended passage than a crossing.

The Farrier does not look up. The Horse does not acknowledge \textsc{Rorario}'s presence in any way that would indicate recognition of him as a socially significant figure requiring response; it registers his passage as a change in the ambient environment and continues its weight-distribution adjustments without modification.

\textsc{Rorario} watches the horse for a duration that is, by a small but perceptible margin, longer than the duration that the passage of the yard strictly requires. The excess is not large --- perhaps three seconds, perhaps four --- but it falls within the interval that separates incidental perception from directed attention, the interval in which something has been looked at rather than merely seen.

He continues through the yard and exits through the far gate.

The Farrier continues working. The horse continues its adjustments.

The sound of the rasp persists briefly in the acoustic field after \textsc{Rorario} has passed beyond the gate's threshold, then is absorbed by the ambient noise of the larger complex.
\end{action}

\transition{CUT TO:}

% ------------------------------------------------------------
\scenenumber{5}
% ------------------------------------------------------------

\sceneheading{EXT. Roads, Post-Moh\'{a}cs Hungary — Day}

\begin{action}
The landscape through which the party moves has sustained the kind of political damage that does not announce itself through dramatic ruin but through a subtler derangement of the normal relationship between territorial markers and the territories they are supposed to designate. Roads that were, within living memory, unambiguously within the jurisdiction of one authority now pass through zones where the nature of that authority has become contested, revised, or simply unclear, the contest having been conducted through a series of military, administrative, and ecclesiastical actions whose cumulative result is a map --- the administrative map, the map of legitimate claim --- that no longer corresponds in any reliable way to the physical territory it purports to describe.

Boundary markers, where they have survived the intervening decades without removal or destruction, designate entities that have themselves been modified, dissolved, or replaced. Some markers have been supplemented by additions in different materials, different scripts, different languages, the supplements contradicting or contextualizing or simply coexisting with the original inscriptions without any mechanism of adjudication between them. Others are entirely absent from their expected positions, the posts still standing but their signifying surfaces missing, removed or damaged so that what remains is the infrastructure of signification without its content --- a column that marks the location of a former boundary marker rather than a boundary.
\end{action}

\begin{action}
\textsc{Rorario} travels as a member of a small party whose composition is functionally adequate for the journey's diplomatic purpose --- enough figures to constitute a credible legation, few enough to move at a pace that the road's current condition permits. The party moves in loose formation, the spacing between its members varying in response to the road's surface variations and the relative positions of those ahead.

Ahead of the party, at a distance of perhaps thirty meters, a \textsc{Rider} makes intermittent progress along the same road. He is, by every available indication, not a member of \textsc{Rorario}'s party but rather an independent traveler whose direction of movement happens to coincide with the party's own. He has halted, for the third time within the period of observation, and is engaged in the consultation of a document --- a map, from the manner in which he holds it deployed before him, its surface oriented upward and outward in the posture of someone cross-referencing a represented space against an experienced one.

The Rider's consultations have not produced, in any of their previous instances, a resolution of the question that motivated them --- or if they have, the resolution has been insufficient to sustain the forward movement of travel beyond the interval between this halt and the next, each halt reimposing the question that the preceding halt had appeared, briefly, to answer. The map is being applied to a territory it cannot fully accommodate; each consultation reveals a new local discrepancy between the representation and the ground, requiring a new act of interpretation that restores provisional progress without establishing the stable correspondence that would make further consultation unnecessary.
\end{action}

\begin{action}
In the same visual field as the halted Rider, sharing the road and its margins without sharing the Rider's difficulty, a pair of \textsc{Oxen} move a loaded cart through a section of road whose surface has been damaged by recent water --- a channel of erosion cutting across the road's width, reducing the passable surface to a narrow band on the road's left margin whose adequacy for the cart's passage is not, from a distance, immediately apparent.

The lead ox encounters the edge of the damaged section and adjusts without pause: a slight leftward pressure against the yoke, a redistribution of pulling force that accommodates the narrowed passage while maintaining the momentum of the forward movement. The adjustment is not a solution to a problem that has been recognized and deliberated upon; it is a continuous modification of the ongoing act of pulling, the road's surface variations registering as inputs into a system whose outputs are the small real-time corrections that keep the cart moving in the direction it is moving. The cart passes through the damaged section. The oxen continue.

The Rider remains halted, the map before him.

Neither the Rider nor the oxen are, within the film's formal logic, positioned as elements of a comparison. They share the road; they occupy the same frame; they proceed, or do not proceed, according to their respective constraint regimes. That these regimes differ in their response latency --- that one produces continuous forward movement and the other produces iterative halts --- is a feature of their co-presence, not an argument constructed from it. The viewer who assembles the comparison does so independently of any editorial direction toward that assembly.
\end{action}

\begin{action}
A \textsc{Dog} joins the party's progress from a lateral direction --- emerging from the margin of the road at a point some twenty meters ahead of \textsc{Rorario}'s current position, falling into a pace alongside the party's forward movement as if the party's trajectory were simply one component of a larger spatial itinerary that the dog has been conducting according to its own organization. It runs parallel for perhaps forty meters with the easy, measured gait of an animal that has not yet reached the exertion threshold that would require a modification of pace, its attention distributed between the road ahead and the party alongside it in a proportion that is difficult to determine from external observation.

Then, without any signal visible to the human members of the party, it veers sharply left and disappears into the vegetation at the road's margin, the departure as unannounced as the arrival, directed toward some object or region whose salience was invisible to everything in the party and apparently absolute to the dog.

The road continues. The Rider, now some distance behind the party, is still consulting the map.
\end{action}

\begin{action}
The landscape's political illegibility does not constitute an obstacle in any straightforward sense --- \textsc{Rorario} and his party continue to move through it, their progress unimpeded by the absence of stable jurisdictional boundaries, which they have long since learned to treat as a background condition of travel in this region rather than as an exceptional circumstance requiring special navigation. The unmarked borders are crossed without ceremony, without the deceleration that marked territorial transitions would occasion, without any moment at which the party can be said to have left one political space and entered another. The transition is accomplished through the simple continuation of movement, its political significance distributed across a spatial interval rather than concentrated at a point.

This is, in its own register, the same operation \textsc{Rorario} performs in every administrative exchange: the crossing of a boundary whose location is unclear by the expedient of continuing to move without requiring the boundary's location to be precisely known before the crossing can be completed. The competence is transferable. The landscape, which has had its maps unmade by the same forces of political reorganization that unmake the institutional frameworks within which \textsc{Rorario} operates, requires the same fundamental adaptive procedure: maintain forward motion; absorb local perturbations; do not require global legibility as a precondition of local action.

The road continues northeast. The sky above it is overcast in the manner that neither promises nor forecloses precipitation, a condition of atmospheric suspension that adds its own low-amplitude underdetermination to the journey's environmental register.
\end{action}

\transition{CUT TO:}

% ------------------------------------------------------------
\scenenumber{6}
% ------------------------------------------------------------

\sceneheading{INT. Secondary Administrative Chamber — Late Afternoon}

\begin{action}
The chamber resembles, in its general proportions and furnishing, the administrative spaces previously encountered, yet differs in ways that resist immediate categorization. Its dimensions are marginally narrower; the placement of its table is offset by a degree that is perceptible without being measurable at a glance; the window, though occupying a similar elevation to that of the writing room, admits light of a different character, filtered through an exterior obstruction whose nature is not visible from within, producing an illumination that is at once sufficient for work and insufficient for the establishment of sharp relational contrasts between objects.

Several individuals occupy the room, though their distribution is less dense than in the antechamber and less formally structured than in the reception chamber. They are engaged in tasks that are recognizably administrative --- sorting, copying, verifying --- yet the coordination between these tasks does not present itself as a unified workflow. Each activity appears locally complete while lacking any immediately legible integration with the others.
\end{action}

\begin{action}
At the central table, two \textsc{Clerks} work in parallel upon documents whose formats are nearly identical. The headings, the spacing of the lines, the placement of seals --- all conform to a shared template that suggests a standardized procedure. Yet upon closer inspection, small divergences become apparent: a variation in phrasing in the second paragraph; a difference in the ordering of clauses; a seal impressed slightly higher on one page than on the other. These divergences do not prevent either document from being legible as an instance of the same administrative form. They do, however, preclude their perfect equivalence.

\textsc{Rorario} enters without announcement. His presence is registered by the nearest Clerk, who inclines his head minimally and returns to his work without interruption. He approaches the table and regards the two documents. He does not immediately take either.
\end{action}

\begin{action}
The first Clerk completes his document, sands it, folds it, and applies a seal with a motion that is both practiced and slightly misaligned --- the seal's impression landing not at the exact center of the intended location but displaced by a margin small enough to fall within the tolerances of acceptability and large enough to be perceptible to an attentive observer. The second Clerk, working from the same template, arrives at the corresponding point approximately six seconds later. His seal is placed with greater precision, yet the wax, having been heated to a slightly higher temperature, spreads more thinly upon contact, producing an impression whose edges are less sharply defined. The two documents now occupy adjacent positions on the table, their similarities more immediately apparent than their differences, yet the differences persist, distributed across their surfaces in a manner that resists immediate synthesis.

A third document lies beneath the two completed ones, partially obscured. Its presence becomes apparent only when \textsc{Rorario} lifts one of the upper sheets slightly, revealing the edge of the underlying page. He draws it out. This document bears the same heading as the other two, the same general structure, yet its content diverges more substantially: a clause present in both upper documents is absent here; another clause appears in its place, addressing a contingency that the other two do not acknowledge. The seal on this document has already been applied. Its wax is fully cooled. \textsc{Rorario} regards the three documents together. The Clerks continue working. Neither looks up.
\end{action}

\begin{action}
A \textsc{Messenger} enters carrying a folded dispatch. He delivers it to one of the Clerks and withdraws without awaiting acknowledgment. The Clerk unfolds the dispatch and reads. His expression does not change. He reaches for a blank sheet and begins to write. The document he produces follows the same general template as the others, yet incorporates a phrasing drawn from the newly delivered dispatch --- a temporal qualifier that was absent from the previous documents, specifying that the instruction it conveys is to be understood as provisional pending further clarification from an authority not named within the document itself. He completes the document, sands it, folds it, and places it upon the table. He does not immediately seal it.

\textsc{Rorario}, without selecting any one document as primary, gathers two of them --- the first completed document and the unsealed one --- and aligns their edges. He does not read them in full. His gaze moves across their surfaces in short, discontinuous intervals, pausing at points where the text diverges, then moving on without reconciling the divergence. He sets them down. He takes up the third --- the one previously sealed --- and holds it for a moment longer than the others. Then he returns it to the table. No instruction is issued. No document is selected for dispatch. The configuration remains.
\end{action}

\begin{action}
The light in the room has shifted subtly during this interval. The obstruction outside the window has altered its relation to the sun, producing a redistribution of brightness across the table's surface. Certain lines of text that were previously legible at a glance now require closer attention; others, previously indistinct, have come into clearer relief. The documents themselves have not changed. Their legibility has.

\textsc{Rorario} withdraws from the table and moves toward the door. As he reaches it, he pauses in the brief interval required to accommodate a minor discrepancy between the expected resistance of the handle and the resistance encountered. The door opens. He exits. The Clerks continue their work. The documents remain on the table, their relations unresolved, their differences neither amplified nor eliminated, persisting as a configuration that has not yet been forced into a decision.
\end{action}

\transition{CUT TO:}

% ------------------------------------------------------------
\scenenumber{6a — Outdated Instructions at the Habsburg Court}
% ------------------------------------------------------------

\sceneheading{INT. Official Chamber, Habsburg Administration — Day}

\begin{action}
The chamber is of greater formal weight than those previously visited --- its furnishings heavier, its surfaces more recently attended to, its architectural proportions communicating the slight upward pressure on dimensional scale that characterizes spaces designed to register institutional authority through the materials and volumes of their construction rather than through any explicit symbolic apparatus. The table is of dark wood whose surface has been polished to a degree that produces reflections of the room's objects in its face, though the reflections are imprecise, the grain of the wood introducing a distortion that renders each reflected object as a slightly softened version of itself.

A \textsc{Habsburg Official} receives \textsc{Rorario} at the table's head. The reception is formal in its procedure and efficient in its execution, the ceremonial elements abbreviated by a mutual understanding that the meeting's substance does not require the full apparatus of its frame. Instructions from Rome are produced --- a sealed document opened with the deliberate attention of someone for whom the act of opening a papal communication carries a significance calibrated to its source.

The document is three weeks old. This is not noted by either party. The Official reads the instructions with the considered attention of someone processing current intelligence, his subsequent responses calibrated to a situation that has, in the interval between the instructions' composition and their receipt, undergone modifications that the instructions themselves cannot reflect. He responds sincerely. \textsc{Rorario} responds to the response sincerely. The exchange is conducted throughout in the register of mutual good faith, each party performing a competent engagement with a reality that neither has reason, within the room's available information, to doubt.
\end{action}

\begin{action}
Outside the chamber's window, a \textsc{Horse} attends to a stone trough. The animal lowers its head, drinks, withdraws, pauses, lowers its head again. The motion is repetitive in the way that biological necessity is repetitive --- the same action executed with the same sufficiency each time it is executed, not because the animal has resolved the question of thirst through accumulated repetition but because thirst is the kind of question that admits resolution only locally and temporarily, requiring the same answer again at the next interval. The drinking therefore has none of the iterative quality of the administrative exchanges in the room --- it does not build toward a conclusion that will render further repetition unnecessary. It simply recurs, each instance complete, none superseding its predecessor, none depending upon it.

The conversation in the room continues. The Official reaches a position. \textsc{Rorario} receives it. He will transmit it. The instructions will have moved one step further through the network before the situation they address has moved even a fraction of the distance that separates it from their description. The Official seals his response. The Horse withdraws from the trough and stands in the mid-afternoon light, its weight settling into an arrangement of uncomplicated equilibrium.
\end{action}

\transition{CUT TO:}

% ------------------------------------------------------------
\scenenumber{7 — Regensburg}
% ------------------------------------------------------------

\sceneheading{INT. Assembly Chamber, Diet of Regensburg — Day}

\begin{action}
The assembly chamber at Regensburg is the largest interior space the film has yet visited, its volume the product of an architectural ambition whose purpose was to make the assembly of many bodies from many jurisdictions feel like a single coherent event. The ceiling is high enough that voices reaching it return altered, their upper registers absorbed by the stone and their lower registers returned as a faint reverberation that underlies all speech with a low-frequency continuous tone, as though the room itself maintains a background of undissolved utterance independent of whatever is currently being said within it.

The Book of Regensburg is being amended. This fact is communicated not by any title page or announcement but by the specific posture of deliberation: the arrangement of bodies around a central text, the presence of multiple hands each bearing instruments for the addition or modification of notation, the particular quality of attention that a document under active revision receives from those responsible for its revision, which differs from the attention of reading in that it is both more focused and more anticipatory, each reader oriented not only toward what the text currently says but toward what it might be made to say and whether that potential saying can be made compatible with the sayings of the others around the table.

The clause under discussion is a single sentence occupying the lower third of one page. Contarini's associates, Melanchthon's interlocutors, Bucer's delegates: each brings to the sentence a framework of interpretive priority that is internally consistent and mutually incompatible with the others in ways that the sentence itself, by virtue of its being a sentence in a natural language whose terms carry multiple implications depending on the theological tradition within which they are read, cannot resolve by becoming more precise. Additional precision in one direction requires the introduction of terms that are themselves contested. The amendment of each contested term opens a new contested term in its vicinity. The document grows without converging.
\end{action}

\begin{action}
At the window's edge, a \textsc{Spider} has been constructing since before the session began. Its web occupies the corner formed by the window's frame and the adjacent stone, a geometry determined entirely by the available anchor points and the properties of the material it produces. The construction proceeds without reference to the proceedings in the room: the spider moves from anchor point to anchor point, extending radial threads and then returning to lay the spiral, each pass adding to a structure whose overall form is a consequence of the method rather than a plan imposed upon the method from outside. By the session's midpoint the web is substantially complete. The spider rests at its center. The clause remains under discussion.

Partially visible behind the leg of a delegate's chair, a \textsc{Dog} has been present since the room filled. It is in the posture of repose, its attention not disengaged but distributed across the room in the continuous low-amplitude monitoring characteristic of animals in occupied spaces. It had a bone at some earlier point --- this can be inferred from its current posture, which is the posture of an animal that has completed something and achieved the particular quality of stillness that follows completion rather than the quality that precedes action. Whether the bone is still present, having been pushed beneath the chair, or has been removed at some point during the session is not determinable from the available visual evidence. The dog sleeps. The clause has not been resolved.

A \textsc{Monk} at the table's middle, listening to a formulation being advanced by one of the delegates, reaches beneath his robe and scratches his left flank with the careful containment of someone performing a private act in public, the gesture minimized but not eliminated, the body's demand accommodated within the constraints of the formal setting by reduction rather than suppression. He returns his hand to the table. The formulation continues.
\end{action}

\begin{action}
\textsc{Rorario} is present in the chamber in a capacity that is formally designated but practically peripheral to the central deliberation. He occupies a position that gives him access to the proceedings without placing him among those responsible for the document's revision. He listens. His listening has, over the course of the session, undergone a qualitative modification that is not immediately legible as such: the attending posture has been maintained, the orientation toward the speakers preserved, yet the quality of his engagement with the content of what is being said has altered in a way that might be described, at the physiological level, as a reduction in the depth of semantic processing --- the sounds and their grammatical structures continue to be received without difficulty, but their elaboration into the fuller cognitive representations that constitute understanding has become, not impaired, but somewhat less automatic, as though the machinery of interpretation were continuing to operate while its motivation had undergone some diminution.

He is not bored. Boredom implies a relationship to time as something to be passed rather than inhabited. What \textsc{Rorario} is experiencing, to the extent that internal states can be attributed to a figure defined primarily by his behavioral signature, is something closer to saturation: the condition of a medium that has received so much of a particular substance that its capacity for further absorption has been temporarily reduced without being eliminated.

Outside, it is afternoon. This is ascertainable through the window above the spider's web, whose light has shifted from its morning angle to the longer, more horizontal approach of a sun declining toward the west. The web, illuminated from this angle, is briefly more visible than it has been at any previous point in the session, each thread catching the light at an angle that renders it perceptible from across the room to anyone whose attention happens to be directed toward the window at this moment.

The clause remains unresolved. The session will continue.
\end{action}

\transition{CUT TO:}

% ------------------------------------------------------------
\scenenumber{7a — Street Theater: The Sacrifice}
% ------------------------------------------------------------

\sceneheading{EXT. Market Square, Regensburg — Afternoon}

\begin{action}
The square outside the assembly building accommodates, in the intervals between its commercial and administrative functions, the kinds of peripheral activity that congregations of people from many places generate: itinerant vendors, small altercations being negotiated, children occupying the margins of adult spatial arrangements, and, at the square's eastern edge, a group of \textsc{Performers} who have established themselves at sufficient distance from the assembly building's entrance to suggest they are not officially connected to the Diet's proceedings while remaining close enough to benefit from its traffic.

The performance is already underway when \textsc{Rorario} passes through the square during a recess. The performers use a simple stage arrangement --- a raised platform of wood approximately knee height, a painted cloth backdrop whose imagery is approximate rather than representational, suggesting landscape and altar through color fields rather than depicted forms --- and their costumes are of the kind that indicates period and role through selected elements rather than comprehensive reconstruction: a particular arrangement of cloth suggesting priestly function, a length of rope suggesting binding, a raised arrangement of wood suggesting the structure upon which something might be placed and prepared.

The figure playing the old man moves with the stiffness that may indicate either the character's age or the performer's, the ambiguity unresolved and unnecessary to resolve. He carries wood. The figure of the young man, already bound, has been arranged upon the raised platform in a posture of resignation that has passed beyond the performed and achieved something close to the actual, the body's arrangement producing in the watching crowd a quality of attention different from that which theatrical representation typically generates.

The knife is raised.
\end{action}

\begin{action}
Several versions of this event have existed in circulation for centuries, their differences neither reconciled nor particularly troubling to those who know them all simultaneously. In the version most familiar from the primary text, intervention arrives at the knife's apex: a voice, a ram, a substitution that redirects the sacrifice before it is consummated, the son surviving to become the ancestor of nations. In other versions, older or later depending on which historiographic framework assigns priority to embarrassment or to authority, the knife descends. The ashes remain. What is subsequently present in the world is not the man who was bound but something that has passed through the altar and returned --- present but altered, the passage having constituted a transformation whose nature the theological traditions that carry this alternative account describe in terms that have their own internal coherence and their external incompatibility with the first account.

The performer playing the angel arrives slightly late --- a timing imprecision that produces a moment in which the knife has begun its descent and the intervention has not yet materialized, an interval of perhaps one and a half seconds in which both outcomes are simultaneously plausible, neither yet foreclosed. The crowd in this interval is still. The quality of its stillness differs from the stillness that precedes a known conclusion.

The angel arrives. The ram is indicated by a second performer emerging from behind the cloth backdrop. The sacrifice is redirected. The crowd's response is of the kind that greets the restoration of the expected, which is not quite the same as relief since the expected was always known to be coming, but which carries some of that relief's somatic character anyway, the body responding to the avoidance of harm regardless of whether the mind knew the harm was never real.

\textsc{Rorario} watches from the square's edge. He has not stopped moving; he has slowed, his transit through the square having assumed the character of an extended passage rather than a crossing. He observes the performance's conclusion. Then he continues through the square and exits at its far side.
\end{action}

\begin{action}
The performers reset for the next iteration. The young man is unbound, repositioned, rebound. The old man retrieves the wood. The angel withdraws behind the backdrop. The audience, partially dispersed and partially reformed from new arrivals, prepares for a performance of a story it already knows, whose interest lies not in surprise but in the quality of execution of the known --- in whether the interval before the intervention will again produce that particular stillness, whether the timing will again be slightly imprecise, and what difference that imprecision makes to the question of which version of the story is being told.
\end{action}

\transition{CUT TO:}

% ------------------------------------------------------------
\scenenumber{8 — The Bubbles}
% ------------------------------------------------------------

\sceneheading{EXT. Village Road — Day}

\begin{action}
The village through which \textsc{Rorario}'s party passes in transit between two points of institutional significance is of the kind that exists in relation to the administrative world primarily as a site of passage and taxation rather than as an object of administrative attention in its own right. Its buildings are arranged along the road's edge with the pragmatic adjacency of structures whose location was determined by access to the road rather than by any principle of spatial planning, each slightly different from its neighbors in the details of its construction while conforming to the same general parameters of material and scale. The road surface is better maintained here than in the more open stretches between settlements, the accumulated interest of inhabitants in the passability of their immediate environment having produced a different quality of surface than the inattention of the open road produces.

A \textsc{Girl} occupies the small open space between two buildings, engaged with a vessel of soaped water and the instrument for producing bubbles from it, a thin loop mounted on a handle through which she draws breath and releases it in the controlled exhalation that the activity requires. The bubbles emerge from the instrument with the particular involuntary beauty of objects whose form is determined entirely by the physics of their production: spherical because that is the form that minimizes surface tension for a given volume, iridescent because that is what thin soap film does to incident light, brief because the conditions that produce them are the same conditions that will destroy them.
\end{action}

\begin{action}
The girl's \textsc{Dog} attends to the bubbles with a quality of concentration that has none of the aesthetic dimension of the girl's own engagement with them --- it is not interested in their form or their optical properties but in their trajectories, which it tracks with the predictive accuracy of an animal whose sensory and motor systems are organized around the detection and interception of moving objects. When a bubble drifts within range, the dog commits: it moves without preliminary, its entire body organized toward the point in space where the bubble will be at the moment of interception rather than where it currently is, the calculation performed instantaneously and below any level that could be described as deliberation. Some interceptions succeed; the bubble is eliminated at the moment of contact, the thin film breaking against the dog's nose or snapping jaws. Some fail; the bubble, altered in its trajectory by the disturbance of the dog's approach, drifts beyond reach and bursts independently against a surface or in the air, the failure accepted without revision of method, the next bubble already being tracked.

The girl does not produce bubbles for the dog. She produces them for herself, for the pleasure of production, and the dog attends to them as a secondary consequence of their existence in the shared space. The bubbles drift without destination. The dog pursues them with total commitment. None are preserved.
\end{action}

\begin{action}
\textsc{Rorario} passes through this configuration at the pace his party sets, which is the pace of horses with a destination and an awareness of the remaining distance. He does not stop. He does not address the girl or the dog. The scene is not staged for his benefit and does not modify itself in response to his passage. He moves through it as through any other feature of the road's environment, his attention touching it briefly and continuing.

He does not look back.

The party continues. The road narrows again beyond the settlement, the buildings withdrawing, the open country reasserting its characteristic quality of spatial indeterminacy.
\end{action}

\transition{CUT TO:}

% ------------------------------------------------------------
\scenenumber{8a — Animal Sequence: Dusk}
% ------------------------------------------------------------

\sceneheading{EXT. Roadside — Dusk}

\begin{action}
The party has halted for the night at a position that is neither settlement nor wilderness, a roadside stop of the kind that accumulates at intervals along traveled routes: a well, a structure for the horses, the remains of fires made by previous travelers visible in a cleared area whose use has been established by repetition rather than by any formal designation.

The horses are attended to first. A \textsc{Groom} moves among them in the fading light, checking legs, removing equipment, distributing water and feed. The horses' responses to his attendance have the quality of animals accustomed to human contact --- not affectionate, not indifferent, but engaged in the practical transaction of their own biological maintenance within the context of human management, each animal alert to its own needs and to the groom's movements in the specific register that those needs and that presence make relevant.

An \textsc{Owl} calls from somewhere in the trees beyond the road's edge. The call has no addressee among the humans present, who register it as ambient sound without response. One of the horses' ears rotates toward it, the rotation precise and instantaneous, a direction-finding adjustment with no sequel. The owl calls again from a different position. The horse's ear rotates again. Between calls, the ear is still.

\textsc{Rorario} sits at the edge of the cleared area, a document on his knees. The light is insufficient for reading. He holds it anyway.

A \textsc{Mouse} crosses the cleared area along the fire circle's near edge, its movement the high-frequency locomotion of a small animal in an open space --- quick, directionally committed, without any of the exploratory circling that larger animals in unfamiliar territory employ. It reaches the fire circle's far edge and disappears into the grass. Its entire transit occupies perhaps four seconds. Nothing in the camp responds to it except one of the horses, who has turned to watch with the mild interest of an animal for whom small fast-moving creatures belong to a category that has some relevance to its attentional system but not, at this moment, enough to produce any further response.

The fire is lit. The light reorganizes the space around it, the near circle rendered warmer and more legible, the surrounding darkness becoming by contrast more absolute. The owl calls again. \textsc{Rorario} sets the document aside.
\end{action}

\transition{CUT TO:}

% ------------------------------------------------------------
\scenenumber{9a — The Farm}
% ------------------------------------------------------------

\sceneheading{EXT. Farmstead — Day}

\begin{action}
The farmstead presents itself as a system already in operation, its rhythms established long before the arrival of \textsc{Rorario}'s party and continuing without modification in response to that arrival. A \textsc{Woman} coordinates the farm's activity with the distributed attention of someone managing multiple concurrent processes in an environment where those processes are not always directly observable simultaneously --- moving between the yard and the outbuildings with a frequency of route-change that suggests continuous prioritization of the most immediate demand, her movement itself a kind of information about which of the farm's several operations is currently furthest from the condition of managed sufficiency.

Animals are distributed throughout: a pair of working oxen in an adjacent field, chickens in a fenced enclosure whose boundary they test continuously without apparent strategic purpose, two dogs at different positions in the yard whose functions appear to be different --- one stationed, attending to the farmstead's perimeter with the focused periodic surveillance of an animal whose role involves monitoring; one mobile, moving in proximity to the Woman's own movements, not herding but attending, close enough to be available for direction without requiring direction to be occupied.

The Woman moves with the economy of someone for whom the physical environment of her work is as well known as the interior of a frequently inhabited room: she does not look at the ground when she crosses it, does not hesitate at the thresholds of outbuildings, does not reassess the position of tools before reaching for them. The knowledge is embodied in the organization of her movement rather than formulated in any register that would require consultation. The farm is running. She is how it runs.
\end{action}

\begin{action}
\textsc{Rorario}'s party has stopped for water and directions. She provides both without interrupting the farm's operation: directing them to the well with a gesture while continuing to manage a situation in the chicken yard whose specifics are not apparent from outside the enclosure but which is absorbing her primary attention. The directions are accurate. The water is where she indicates. She does not come to the party; she receives it at the margin of her actual activity.

At the far edge of the yard, one of the oxen has halted in the adjacent field and is engaged with some feature of the ground that is not visible from the yard --- perhaps an irregularity in the surface, perhaps nothing that would be visible to any observer, the animal's attention directed toward a distinction in the ground that exists within a sensory register not available to human observation. It stands there for a duration that exceeds what task completion would require, then moves on without indication of what it found or failed to find.

Something physical and slightly ungovernable is underway at the yard's edge near the pig enclosure --- a smell of considerable directness, a sound of biological process completing itself without reference to the aesthetic preferences of those nearby, the animal world's fundamental indifference to the categories of propriety proceeding with its characteristic thoroughness. No one reacts. The farm continues.

\textsc{Rorario}'s party departs. The woman has not paused.
\end{action}

\transition{CUT TO:}

% ------------------------------------------------------------
\scenenumber{9 — The Bridge}
% ------------------------------------------------------------

\sceneheading{EXT. Road and Bridge — Day}

\begin{action}
The bridge presents itself at a point in the road's descent toward a river crossing, its stone construction visible from a sufficient distance that the party has time to assess it before arrival. The assessment, conducted without discussion among the riders, produces a collective deceleration as the party approaches --- not a halt, but a reduction in pace that indicates a shared uncertainty regarding the crossing's advisability, an uncertainty produced not by any specific visible defect but by the aggregate of small indicators that experienced travelers learn to read in bridge construction: the waterline staining on the piers, the quality of the mortar at the joints, the slight unevenness of the roadbed surface approaching the near arch.

The \textsc{Lead Horse} stops at the bridge's approach, not suddenly but with the gradual deceleration of an animal that has decided, through whatever the perceptual and evaluative processes available to it produce, that forward movement into this particular space is not the action it is prepared to execute. It does not back away. It stands. Its rider applies pressure. The horse receives the pressure, incorporates it into the balance of its current state, and does not move. The pressure increases. The horse's head moves slightly, an acknowledgment of the communication without constituting compliance with its direction.

The party crosses by a lower ford some four hundred meters upstream, a route longer and less comfortable than the bridge would have been, the horses picking their way through the shallow water with the care that uneven footing requires. The bridge remains behind them.
\end{action}

\begin{action}
Later --- not in the next scene but at a point in the film's temporal extension that can be measured only by the accumulation of scenes between this one and that --- a sound will be heard from a distance: the particular acoustic signature of a stone structure under sudden vertical load whose distribution the structure cannot sustain. The sound is brief. It is heard by whoever is nearby at that moment, and by no one with a specific connection to the party's earlier crossing by the ford. The bridge is not mentioned again. The horse is not mentioned again. These events do not form a narrative unit. They are simply events that occurred in proximity, separated by time, connected by nothing that the film acknowledges as connection.
\end{action}

\transition{CUT TO:}

% ------------------------------------------------------------
\scenenumber{10a — The Falcon}
% ------------------------------------------------------------

\sceneheading{EXT. Hunting Grounds — Afternoon}

\begin{action}
The return from the hunt. The \textsc{Falcon} is on the glove of its \textsc{Handler}, hooded, the hood performing its function of reducing the visual field to zero and thereby reducing the animal's arousal state to the level at which it can be carried without requiring the continuous restraint that full visual engagement would produce. The Handler speaks to it. He has been speaking to it throughout the return, a continuous low-register address whose content mixes functional information, assessment of the hunt's outcomes, and a kind of habituated commentary whose purpose is less communicative than tonal --- the voice as a carrier medium for a quality of presence that the animal attends to without processing at the level of semantic content.

The falcon's responses to this speech are not responses to its meaning. They are responses to its prosodic features: changes in volume, changes in rhythm, the onset and cessation of utterance. When the Handler's voice rises slightly, the falcon's grip tightens on the glove. When his pace changes, the bird's body adjusts to maintain its balance relative to the new motion. The adjustments are immediate, executed within the response interval of a sensory-motor system calibrated to handle much faster events than human speech, and they track the physical information embedded in the speech while remaining entirely independent of its propositional content.

The Handler is not wrong to speak as he speaks. The speech is doing something. It is simply not doing what speech between humans does, and this difference is not visible in the speaking, only in the attending.

\textsc{Rorario} rides within hearing distance of the Handler's continuing address. He does not look toward them. The afternoon light produces long shadows from the returning party's figures and horses, their extended dark forms preceding them across the field.
\end{action}

\transition{CUT TO:}

% ------------------------------------------------------------
\scenenumber{11 — The Registry Hall}
% ------------------------------------------------------------

\sceneheading{INT. Ecclesiastical Registry Hall — Day}

\begin{action}
The hall is larger than any interior previously encountered, its volume subdivided by rows of tall shelving structures whose vertical extension exceeds that of a single story, producing the impression of a space that has been internally stratified rather than architecturally expanded. The shelves are densely occupied by bound volumes, folded documents, and loose fascicles inserted between more stable forms, their arrangement governed by a classificatory system that is at once highly articulated and only partially legible from within the space itself.

Light enters from a series of high clerestory openings, descending in oblique shafts that illuminate sections of the shelving while leaving adjacent regions in relative obscurity. The effect is not one of dramatic contrast but of uneven legibility, certain records available to immediate perception while others remain present but unarticulated within the field. Movement through the hall is continuous but non-uniform: Clerks traverse the aisles along paths that appear habitual rather than prescribed, each adjusting minimally to the presence of others in a manner that suggests a shared but unformalized spatial logic.

\textsc{Rorario} enters accompanied by a \textsc{Senior Clerk} whose familiarity with the space is evident in the absence of hesitation. They arrive at a table positioned within one of the illuminated zones. Two volumes are already present upon it.
\end{action}

\begin{action}
The volumes are of comparable size and binding, their covers bearing inscriptions in similar hands, yet the differences between them become apparent upon closer inspection: the coloration of the ink, the slight variation in the thickness of the binding, the degree of wear along the edges. Each presents itself as an authoritative record within the same administrative domain.

The Senior Clerk opens the first volume. Its pages contain entries arranged in tabular format, each line recording a determination regarding jurisdictional authority over a series of contested regions, proceeding in chronological order, building upon the prior, establishing a sequence of decisions that together constitute a coherent administrative history. He opens the second volume. Its structure mirrors that of the first --- tabular entries, chronological progression, formal consistency --- yet the determinations it records diverge. Where the first volume assigns authority to one office, the second assigns it to another; where the first records a resolution, the second records a deferral; where the first indicates closure, the second preserves openness. The divergence is not random. It is systematic.
\end{action}

\says{Senior Clerk}{Both registers are in current use.}

\begin{action}
The statement is delivered without emphasis, without qualification. It does not present itself as a problem. It presents itself as a condition.

At an adjacent table, another \textsc{Clerk} copies from a third volume into a fresh document, his transcription proceeding without interruption. The content he copies aligns, at intervals, with the first volume; at other intervals, with the second; at still others, with neither, incorporating phrasing that does not appear in either register but conforms to their general format. He does not pause to reconcile these differences. He writes.

A \textsc{Messenger} approaches, carrying a sealed instruction. The Senior Clerk breaks the seal and reads, his expression unchanged, then turns to \textsc{Rorario}.
\end{action}

\says{Senior Clerk}{The determination is to be enforced immediately.}

\begin{action}
He does not specify which determination. He does not indicate which volume is to be taken as authoritative. He hands the instruction to \textsc{Rorario}, who reads. The instruction refers to the matter recorded in both volumes. It does not reference either volume explicitly. Its phrasing is compatible with both determinations. \textsc{Rorario} folds the instruction and places it upon the table between the two volumes. The three objects --- the first register, the second register, the instruction --- form a configuration whose relations are fully specified at the level of physical arrangement and entirely unspecified at the level of operational hierarchy.

At the hall's far end, two groups of Clerks prepare documents for dispatch, each drawing from different sections of the shelving, each producing documents bearing the seals of distinct offices. The documents are bundled. Messengers receive them. They depart through different exits. The effect is not one of contradiction. It is one of parallelism: two sequences of action, each internally coherent, each proceeding from an authoritative source, extending outward into the institutional network without intersecting within the space of their production.

The Senior Clerk closes both volumes. The sound of their covers meeting their pages is nearly identical. He returns them to their respective shelves. They are placed in different sections. The instruction remains on the table. A Messenger, not the one who delivered it, approaches, takes it without inquiry, and departs. No record is made of which determination it will enact.
\end{action}

\transition{CUT TO:}

% ------------------------------------------------------------
\scenenumber{12 — The Boundary}
% ------------------------------------------------------------

\sceneheading{EXT. Boundary Zone — Early Evening}

\begin{action}
The boundary is not marked by a line. It is indicated, if at all, by a gradual change in the condition of the land: a shift in the density of vegetation, a variation in the maintenance of the road, the appearance of structures whose architectural features correspond more closely to one administrative region than another. These indicators do not converge at a single point; they overlap across a spatial interval within which the designation of the boundary is a matter of interpretation rather than observation.

A small group has assembled. \textsc{Rorario} stands at its center, the document in his possession. Opposite him, a \textsc{Local Official} represents the authority whose jurisdiction is to be affirmed or altered. Several others are present --- assistants, observers, figures whose roles are not explicitly declared but whose presence constitutes an audience sufficient to register the act as official. A second official, positioned slightly behind the first, advances an alternative reading of the instruction's designation of the boundary's location. Both statements coexist within the same spatial and temporal frame. Neither is withdrawn. Neither is prioritized.

\textsc{Rorario} reads the document again. Its text remains unchanged. The light, now lower in the sky, renders the page more difficult to read. He lowers the document. He looks at the ground. The surface bears no mark that would resolve the question. The group waits.
\end{action}

\begin{action}
\textsc{Rorario} steps forward. He indicates a point on the ground. The gesture is precise. It is not accompanied by explanation. The Local Official nods. The second official does not object. Assistants begin to act. A marker is placed. The action proceeds.

At a short distance, another group --- previously outside the frame of attention --- has begun a similar process. Their marker is placed at a different location. The distance between the two markers is small but non-zero. Both processes continue.

A \textsc{Dog} moves along the interval between the markers. Its path intersects the projections of both groups' measurements, crossing lines that have been established without altering its trajectory in response. It pauses briefly at one of the secondary indicators, attending to a feature of the ground not visible within the human framework of measurement. It moves on.

The light continues to diminish. The boundary, such as it is, now exists in two locations. Neither is erased. Neither is confirmed as singular. The markers remain. Their relations have changed. They have not been resolved.
\end{action}

\transition{CUT TO:}

% ------------------------------------------------------------
\scenenumber{12a — The Market Stall}
% ------------------------------------------------------------

\sceneheading{INT. Market Stall — Late Morning}

\begin{action}
The stall is narrow, its interior partially enclosed by wooden panels whose placement creates a threshold between the open movement of the market and the more contained space within, the boundary between these zones enforced less by the panels' physical opacity than by the conventional understanding that the stall's interior is the Seller's domain and the exterior the Buyer's, a distinction that commerce maintains without requiring it to be declared. Various goods are arranged along the stall's surfaces --- tools, small containers, folded materials --- each placed with a degree of order that suggests recent handling without having achieved a stable configuration, the arrangement partway between organization and the intermediate state that precedes it.

A \textsc{Buyer} stands at the entrance. A \textsc{Seller} is present within. The Buyer has come because he was told, at some prior point, that a particular item would be ready for collection by this hour. Whether the Seller was the source of this information, or whether it arrived through an intermediary, is not established by the scene.
\end{action}

\begin{action}
The Seller is occupied. He holds a small collection of tokens or coins arranged in his hand, which he counts --- though the counting does not proceed in a linear sequence but in a series of partial progressions, each advancing to a point at which some condition causes it to restart, the condition not visible from outside the Seller's own cognitive process. He rearranges the objects in his hand. He begins again. The Buyer waits, which is to say he performs the physical arrangement of waiting --- the stillness of the body, the orientation toward the point of expected activity --- without the waiting producing any result within the scene's available duration.
\end{action}

\says{Buyer}{I was told this was ready.}

\begin{action}
The Seller does not look up immediately. He completes a partial sequence of the counting, stops before its conclusion, and sets the objects down on the stall's surface. He turns to a shelf behind him and retrieves a wrapped item. He examines the wrapping with the attention of someone assessing the adequacy of a prior act, adjusts it slightly, then sets the item down again without presenting it to the Buyer. The interval between retrieving and setting down is occupied by an assessment whose criteria are not communicated.
\end{action}

\says{Seller}{It is nearly so.}

\says{Buyer}{May I see it.}

\begin{action}
The Seller receives this not as a refusal but as a step requiring preparation --- a reasonable request whose satisfaction requires that certain prior conditions be met before the presentation can properly occur. He unwraps a portion of the item, the wrapping partially opening to reveal the object within, then pauses, the act of unwrapping having apparently introduced a new condition whose satisfaction takes precedence over the completion of the unwrapping. He rewraps it.
\end{action}

\says{Seller}{It must be accounted for first.}

\begin{action}
He returns to the tokens. He counts again. The sequence differs from the prior one, proceeding further in one direction before encountering whatever condition causes the restart, but the difference in progress does not translate into a different outcome: the counting again stops before producing a number that could function as its conclusion. The ledger is consulted. Pages are turned; a finger traces a line; a page is turned back. The ledger is closed.
\end{action}

\says{Seller}{There is a discrepancy.}

\begin{action}
The statement is delivered without concern. It does not specify the nature, magnitude, or location of the discrepancy. It does not halt the process; it names a condition within the process, a condition whose resolution is implicit in the Seller's continued engagement with the tokens and the ledger, though this resolution does not appear within the scene's duration.
\end{action}

\says{Seller}{It will be resolved shortly.}

\begin{action}
The Buyer waits. Outside the stall's entrance, the market continues at its ambient pace --- a \textsc{Dog} passes across the threshold opening without pausing, its trajectory carrying it through the scene's frame without reference to what the frame contains. A \textsc{Second Buyer} approaches from behind the first, observes the configuration of the stall's activity, and withdraws without speaking, having determined from what is visible that the conditions for a transaction are not currently met. The first Buyer does not turn to observe this. The Seller does not register it.

The Seller places the ledger beside the wrapped object. He does not open the object again. He returns to the tokens. The light in the stall shifts as the sun's position relative to the panels changes, rendering certain objects on the shelves more legible while others recede into comparative obscurity. This change in the distribution of legibility does not alter the objects themselves. The item remains wrapped. The account remains partially unresolved. The Buyer remains at the threshold. The process continues in a condition of perpetual near-completion that does not complete within any interval the scene makes available.
\end{action}

\transition{CUT TO:}

\sceneheading{INT. Clerical Refectorium — Evening}

\begin{action}
The refectorium is modest in scale, its architecture subordinated to function without the complete reduction of ornament that characterizes purely administrative spaces. A long table occupies the room's center, its surface bearing the remains of a meal in progress: bread partially consumed, cups at varying levels, a shared dish from which portions have been taken without formal division. The lighting is low and uneven, provided by candles whose distribution produces localized zones of clarity separated by regions of softer obscurity.

Several \textsc{Clerics} are seated. \textsc{Rorario} is present among them, not as a primary participant but as one whose presence is accommodated within the existing arrangement. The conversation has already been underway for some time. It does not begin at its beginning.
\end{action}

\says{First Cleric}{It is clearly stated that David struck him down.}

\says{Second Cleric}{It is equally stated that Elhanan did so.}

\begin{action}
The statements are delivered without escalation. They do not initiate the contradiction. They resume it, as one resumes a task that has been set aside not because it was completed but because the available time for it had been exhausted and has now been restored.

The textual situation they are navigating is one of those genuinely irresolvable discrepancies that the scriptural corpus carries within itself without distress: the first account assigns the killing of the Philistine champion to the shepherd-king whose career it inaugurates; the second assigns the same killing --- or a killing indistinguishable from it within the narrative's terms --- to one of David's fighting men, a figure of lesser narrative prominence whose claim to the deed is, if the criterion of authenticity is embarrassment rather than authority, arguably the stronger one, the story that serves no dynastic purpose having less reason to have been introduced.
\end{action}

\says{Third Cleric}{The later recension specifies that it was the brother of Goliath.}

\begin{action}
A brief interval follows. The statement is not rejected. It is not adopted. It is the kind of scribal adjustment --- the insertion of \textit{brother of} into a verse that previously had no such qualification --- that resolves a contradiction at the level of the text while leaving intact the question of what it resolved and whether resolution of that kind constitutes the same thing as truth. The emendation is detectable; it has the slightly forced quality of a joint that has been reinforced after the fact, its reinforcement visible in the altered material around it.
\end{action}

\says{First Cleric}{That is not the earlier text.}

\says{Second Cleric}{It is the text we use.}

\begin{action}
The distinction between \textit{earlier} and \textit{used} is introduced and allowed to remain. No hierarchy is assigned between these criteria. A manuscript lies open on the table between them, its script clear, its lines evenly spaced. One passage is marked. Another passage, on a different page, is also marked. The marks do not correspond to one another in any system that the room contains. A \textsc{Scribe} at the table's end copies a passage from the manuscript into a smaller book. He pauses briefly at the point where the contested name appears. He writes. He does not alter it. He does not annotate it. He continues.

The conversation disperses into adjacent topics without transition. The names are not mentioned again. The manuscript remains open. The marks remain where they were placed. They are not reconciled. \textsc{Rorario} rises without referencing the discussion. He exits. The others remain. The meal continues.
\end{action}

\transition{CUT TO:}

% ------------------------------------------------------------
\scenenumber{13b — The Gateway Mark}
% ------------------------------------------------------------

\sceneheading{EXT. City Gateway — Late Afternoon}

\begin{action}
A stone gateway marks the transition between two districts of the city, though the distinction between the districts is not visibly encoded in the surrounding terrain --- their difference is administrative rather than physical, a designation maintained through records and customs rather than through any difference in the character of the streets or buildings on either side. Traffic passes through without interruption, the threshold crossed many times each hour by those for whom it constitutes nothing more than a passage through a stone arch.

On the interior face of the arch, a mark has appeared. Its form is not standardized against any legible system of notation or symbol: it resembles neither the administrative marks applied to walls and surfaces by institutional authority nor the personal marks of individuals claiming or disclaiming ownership. It is not random; it exhibits a coherence of form that suggests intention, but the intention is not interpretable from the form alone. It has been there long enough to have dried fully into the surface, but its color retains sufficient freshness to indicate that it did not predate the current season.
\end{action}

\begin{action}
A passerby notices it. He pauses briefly, touches its surface with a fingertip in the gesture of someone verifying a material property --- the texture, the dryness, whether it transfers --- and continues. Another figure moving through the gateway avoids looking at it entirely, the avoidance having the character of a deliberate redirection rather than simple inattention. A third traces its outline with his gaze without approaching.

Two individuals confer quietly in the gateway's shadow.
\end{action}

\says{First Passerby}{It signifies protection.}

\says{Second Passerby}{It signifies warning.}

\begin{action}
They do not reconcile their interpretations. They proceed through the gate in the same direction, the mark behind them. A guard stationed at the gateway's margin does not acknowledge the mark; he regulates passage as before, his attention directed toward the flow of people rather than toward the arch's surface. A child passing through the gateway in the company of an adult pauses and attempts to reproduce the mark on an adjacent surface using a stone fragment. The result differs from the original in several of its properties. No correction is offered. The original mark remains, its presence modifying the behavior of those who notice it without establishing a shared account of what it means or requires.
\end{action}

\transition{CUT TO:}

\sceneheading{INT. Negotiation Chamber — Day}

\begin{action}
A long negotiation. The chamber is of intermediate scale, its furnishings adequate to the meeting's duration without being designed for extended occupation: chairs of moderate comfort, a table whose surface bears the accumulated evidence of a session that has been in progress for some time --- a cup, displaced documents, a pen set down and not retrieved.

\textsc{Rorario} is present and responds when addressed. The negotiation's subject has been established prior to the scene's entry, its terms familiar to all parties from prior sessions. He has participated in this meeting's earlier phases, his contributions audible in the record of positions taken and language proposed. He is present now in the same formal capacity.

His contributions have become shorter.

This is not a dramatic change. It is a change of amplitude. Where he would previously have developed a position through a sequence of clauses, each adding a qualification or an extension that made the overall formulation more precisely calibrated to the available space between the parties' interests, he now tends toward the same position in fewer clauses, the compression not producing imprecision but not producing the additional precision that the additional clauses would have introduced. The formulations are sufficient. They are also minimal.

The others do not notice. The negotiation continues to receive his contributions at their new amplitude without registering the reduction as a change in kind, because the change has occurred too gradually to establish a baseline against which it can be measured from within the session.
\end{action}

\begin{action}
All three pressure tracks are simultaneously present and jointly irreconcilable. The geopolitical instability of the territories under discussion has introduced into the session's prior phases a series of contextual modifications each of which required adjustment of the positions being negotiated, such that the positions currently on the table bear a complex genealogical relationship to the positions with which the session began, connected by a chain of accommodations no single one of which was decisive but whose accumulation has produced a significant displacement. The institutional sources from which \textsc{Rorario} draws his mandate have, during the course of the negotiations, issued supplementary instructions whose relationship to the original mandate is one of partial modification, the nature and extent of which modification is not entirely clear from the instructions' text. The animal world has not entered this room. No animal has been in this room.

The negotiation continues. His next contribution, when it comes, is two clauses. The one after that is one clause. The one after that is a single affirmation without elaboration --- a gesture of continued participation rather than a contribution of substance, though it is received as the latter and has the functional effect of the latter, which is to say: the negotiation continues. He is still present. He continues to function. The pop has occurred and no one is aware of it, including, in any register available to behavioral observation, \textsc{Rorario} himself.
\end{action}

\transition{CUT TO:}

% ------------------------------------------------------------
\scenenumber{14a — Regional Chamber}
% ------------------------------------------------------------

\sceneheading{INT. Regional Chamber — Evening}

\begin{action}
The chamber is larger than those previously encountered, though not proportionally more ordered. Multiple tables occupy the space, arranged according to no single geometry but rather by successive accommodations made over time --- each placement preserving a prior necessity while introducing a minor deviation into the overall configuration, the cumulative effect being a room whose furniture arrangement encodes its own administrative history in a form not immediately legible as such but detectable in the slight awkwardness of each table's relation to its neighbors. Documents are distributed across the surfaces in layered strata, older sheets beneath and more recent ones above, though the distinction is not always legible without examination of the paper's condition and the ink's degree of oxidation. Several Clerks are present. A Messenger stands near the entrance holding sealed correspondence that has not yet been received into the system of recording and therefore exists, at this moment, in a condition of administrative suspension --- present in the room, not yet present in the record.

At the central table, \textsc{Rorario} stands with a document in hand. A \textsc{Registrar} indicates a passage already marked.
\end{action}

\says{Registrar}{This reflects the established condition.}

\begin{action}
He gestures to a clause. He produces a second document. The wording differs. The structure does not. A third document is introduced by another Clerk, positioned such that all three may be read simultaneously, though not without shifting the angle of attention required to maintain them in view. Each is internally consistent. Their relation is not resolved.
\end{action}

\says{Registrar}{The revision has been noted separately.}

\says{Clerk}{This corresponds to the recognized form.}

\begin{action}
At a side table, two Clerks compare entries, reading aloud in alternating sequence as though the alternation itself constitutes a method of comparison rather than merely a convention of shared reading.
\end{action}

\says{Clerk One}{As previously established ---}

\says{Clerk Two}{As subsequently recognized ---}

\says{Clerk One}{As locally applied ---}

\says{Clerk Two}{As provisionally amended ---}

\begin{action}
The phrases accumulate. No single phrase invalidates the others. Their joint application cannot be fully determined because the conditions that would allow determination --- a stable hierarchy among the cited authorities, a temporal sequence sufficient to establish precedence, a spatial designation clear enough to assign local applicability --- are each present in partial form and collectively insufficient.

A Messenger approaches with a sealed letter. It is opened. The contents are read silently. The Registrar nods, placing the letter atop one of the documents. The alignment is approximate. The seal is retained, set aside, not discarded, as though its preservation remains relevant even after the letter has been opened and received --- the seal as evidence of a prior condition of closure that the opening has ended but not annulled.

At the far end of the chamber, a small group confers over a diagram --- lines intersecting, diverging, recombining. The diagram corresponds imperfectly to the maps previously observed in the film, sharing their general notational conventions while representing a territory at a slightly different scale and under slightly different administrative assumptions, the discrepancy between the two representations not resolvable within either without reference to the territory itself, which is not present.
\end{action}

\begin{action}
\textsc{Rorario} turns slightly, observing the room as a whole. No single vantage allows all operations to be apprehended simultaneously. A Clerk approaches with a document requiring acknowledgment.
\end{action}

\says{Clerk}{It must be signed to be entered.}

\begin{action}
\textsc{Rorario} examines the document. The clauses refer to prior conditions not fully present in the available material.
\end{action}

\says{Rorario}{It refers to what has not been settled.}

\says{Clerk}{It refers to what has been recorded.}

\begin{action}
A brief interval. \textsc{Rorario} signs. The document is entered. Its inclusion alters the relation among the existing set without altering any element within that set, the alteration occurring at the level of the configuration rather than the level of any constituent.

At another table, a Clerk attempts to reconcile two entries by drawing a line between their corresponding passages, a graphic act of proposed alignment that has no formal status within the room's documentary system and will not survive into any subsequent iteration of the documents it marks. He pauses, then leaves it. A candle is replaced; the light shifts, changing the legibility of certain passages while obscuring others. No adjustment is made to compensate. The Messenger, now carrying additional sealed letters, remains. His continued presence is no longer remarked upon. The chamber continues, each operation completing locally, the system as a whole not converging. It persists.
\end{action}

\transition{CUT TO:}

\sceneheading{INT. Writing Room — Night}

\begin{action}
\textsc{Rorario} alone. A candle. A writing table. A document in progress.

The room is the same species of functional enclosure as all the writing spaces he has occupied throughout the preceding scenes, distinguished from them primarily by the hour and its consequence for the light: a single candle's illumination, directed downward onto the page, leaving the room's peripheral regions in a darkness that is not absolute but is sufficiently deep that the room's dimensions are no longer determinable without movement through it.

He writes. The pen moves across the page with the practiced continuity of many years of similar activity, producing the administrative script that has been the primary medium of his professional existence, its letterforms efficient and legible, their production automatic at the level of execution while remaining deliberate at the level of content.

He writes a sentence. The sentence has begun as such sentences begin, with the formulaic opening that identifies the document's status within the institutional network --- its origin, its addressee, its relation to prior communications. The sentence extends through its necessary subordinate clauses, accumulating the qualifications that precision requires. It reaches the point at which the sentence's central claim must be made, the point for which all the preceding structure has been preparation.

The pen stops.

Not because the ink has failed or the hand has cramped. The pen stops because the continuation has not arrived. The sentence exists up to its central claim in a condition of structural preparedness for a statement that does not come. The pen remains in contact with the page, its nib leaving a small additional deposit of ink at the point of arrest that distinguishes this stop from a deliberate pause, marking it as something that was not planned.

\textsc{Rorario} holds the pen. He does not look at the page. He does not look away from it. His attention is not elsewhere; it is simply not doing what it was doing.

The candle moves slightly --- a draft from somewhere in the building's infrastructure, some gap in the wall's mass --- and the page's illumination shifts. The incomplete sentence is briefly more, then less legible.

He sets the pen down. He gathers the document. He folds it along lines that are not the lines it would have been folded along had it been complete. He seals it. The seal is applied. The letter is a completed object in every sense except the sense in which it contains what a letter is supposed to contain.

He sets it on the pile of completed dispatches.

It will be sent.

The candle burns. End of State~II.
\end{action}

\transition{CUT TO:}

% ------------------------------------------------------------
\scenenumber{15a — Two Altars}
% ------------------------------------------------------------

\sceneheading{EXT. Rural Clearing — Late Afternoon}

\begin{action}
The clearing is unmarked, its boundaries defined only by the thinning of trees at its edges and the condition of the ground, which bears traces of intermittent use without any formal designation of purpose. Two low stone structures occupy the space, positioned at a slight remove from one another, their construction similar but not identical: one composed of more regular stones, its surface relatively even; the other assembled from irregular fragments, its form less stable, its edges uneven. Upon each, an offering has been placed. On the first: cut vegetation, arranged with a degree of care that suggests prior handling, a precision in the layering that implies attention to form as well as function. On the second: a small animal, recently killed, its body positioned without evident ceremony, the placement pragmatic rather than arranged.

No figures are immediately present. The light is low, the shadows elongated, the surfaces of the two structures receiving illumination at slightly different angles due to their relative orientations within the clearing.
\end{action}

\begin{action}
Birds arrive. At first a single bird lands upon the edge of the structure bearing the vegetation. It pauses, its head moving in small increments as it surveys the surface. Then a second. Then several more. They gather without apparent coordination, their distribution across the surface forming a pattern that is neither symmetrical nor random but responsive to local conditions not fully legible from a distance --- the texture of the stone, the arrangement of the vegetation, the particular quality of the afternoon light at this angle of incidence. None of this is visible as cause from outside. The birds are simply there, attending to whatever has drawn them there, which is their business and not the scene's to disclose.

On the second structure, the animal remains undisturbed. A fly lands. Then another. The activity there is of a different order: smaller, less visible at a distance, more localized to the surface of the offering itself, which is to say it is following the biological logic of decomposition rather than the behavioral logic of the birds, and the two processes, though occurring within the same spatial field, have no relation to each other except that of temporal proximity.
\end{action}

\begin{action}
Wind passes through the clearing. The vegetation shifts slightly. A portion of it slips from its arrangement, falling to the ground. The birds adjust. Some depart. Others remain. The structural consequence of this redistribution cannot be evaluated from the available information.

At the edge of the clearing, two \textsc{Figures} become visible. They do not approach immediately. They observe. One gestures, indicating the structure with the animal. The other looks toward the structure with the birds. No words are spoken. The gesture is not repeated. The observation does not resolve into agreement or into its absence.

A brief disturbance: one of the birds takes flight abruptly, followed by several others, the movement propagating across the group in a pattern that suggests response to a stimulus not visible within the frame. Within seconds the structure is nearly empty. One bird remains. Then it too departs. The clearing returns to stillness. The two structures remain. The offerings remain. Their condition has changed. Their relation has not been determined. The Figures turn and leave. They do not approach either structure. They do not alter them.
\end{action}

\transition{CUT TO:}

% ------------------------------------------------------------
\scenenumber{15b — The Spewing Plant}
% ------------------------------------------------------------

\sceneheading{EXT. City Wall, Eastern Quarter — Morning}

\begin{action}
At the base of the city wall, in a narrow interval between the wall's foundation stones and the beginning of the adjacent building's footprint, a plant has established itself in conditions that its establishment renders implausible: the soil is minimal, the light intermittent, the drainage aggressive. It is not a beautiful plant. Its stems are succulent and irregular, its leaves broader than their support would seem to warrant, its overall aspect one of exuberant biological overreach relative to the resources available to it --- a system that has committed more structure than the environment can reliably sustain.

Several people pass along the street without attending to it. A \textsc{Child} squats briefly to examine it with the close attention that children apply to things at ground level that adults have learned to overlook, then moves on. A \textsc{Dog} sniffs at the base of the wall nearby, finds nothing that requires further investigation, continues along the wall's face.

By midday the plant's most extended leaf has begun to flag, its cellular pressure insufficient to maintain the turgor that had held it extended. By early afternoon the stem droops. The process is not dramatic --- it has the quality of a system withdrawing an excess commitment it had never been able to sustain indefinitely, the return of ambition to equilibrium, which is not death so much as the correction of an overclaim.

By evening the plant is substantially reduced from its morning condition. A \textsc{Passerby} who had walked this street in the morning and walks it again now would find, if they had attended to it then, that what had been a considerable presence is now a considerably lesser one, and might form the impression that something had been taken from it, though nothing has been taken --- it has simply become less than it briefly was, and the world has not adjusted itself to accommodate the loss. The wall remains. The gap in the pavement remains. The plant remains, diminished. Tomorrow it may recover. It may not.
\end{action}

\transition{CUT TO:}

% ============================================================
\statebreak{III}{Writing}
% ============================================================

\noindent{\itshape Language no longer functions as a tool of negotiation. It has become the residue of eliminated possibilities. The environment outside continues indifferently. Perceptual flickers are now longer and not always restored. Animals are no longer peripheral to the human scenes but occupy the same visual grammar, the same spatial register, without announcement. The study is the primary space. Rorario is no longer composing. He is at the limit of what the prior system can produce.}

\vspace{1.5em}

% ------------------------------------------------------------
\scenenumber{16 — The Manuscript}
% ------------------------------------------------------------

\sceneheading{INT. Study — Day}

\begin{action}
A room quieter than anything previously occupied in this film, its quiet a function not of soundproofing or distance from activity but of a reduction in the number of active channels that the space is organized to support. Where the antechamber had accommodated four or five simultaneous communicative threads, where the assembly hall at Regensburg had sustained a dozen, this room sustains one: the relation between a figure and a surface, both of which are present and neither of which is, at this moment, adding to the other.

The manuscript exists on the table in a state of advanced but not final form. This is evident in the physical character of the pages: some are clearly settled, their ink fully dried, their surface bearing the slight sheen of pages that have been handled enough times to have absorbed the oils of that handling into their fibers; others are less certain, their placement among the gathered pages tentative in the way of material not yet sure of its position within the whole. The pages have been accumulated rather than produced in a single sustained session --- they bear the marks of different inks, slightly different pressures, the minor variations in letterform that accumulate across time even within the same hand. They represent the sedimented output of a process whose duration corresponds to most of the preceding film.
\end{action}

\begin{action}
\textsc{Rorario} handles the pages without reading them in the sense that requires the semantic processing of content. His engagement with them is more nearly tactile --- he lifts one, turns it slightly, sets it down in a position that may or may not be the position from which he lifted it. He picks up another. He holds it without bringing it to a reading angle. The relationship between the person and the pages he has produced has shifted from the productive to something closer to the curatorial, an attention to the physical existence of the thing rather than to its content as such.

A \textsc{Courier} passes the window on a horse. The horse and its rider cross the window's frame in a duration of perhaps two seconds, moving at a pace that indicates purpose and a destination external to the study. The sound of the hooves on the courtyard surface arrives at the window slightly delayed from the visual, the acoustic lag small but present, the world outside the window governed by its own temporal logic without reference to the study's. The courier does not slow. The study does not acknowledge the passage.

A \textsc{Cat} enters the room. It has been in the building already --- the way it crosses the threshold has the quality of a system re-entering a familiar space rather than entering an unfamiliar one, its movement confident in the sense of not requiring reconnaissance before commitment to a path. It crosses the study along a trajectory that appears, from the available evidence, to be directed toward a specific destination in the room's far corner, though what draws it there is not apparent. It arrives. It settles into the brief investigatory stillness that precedes a cat's decision about whether a location merits occupation. It decides against. It leaves by the same door through which it entered. The manuscript remains. \textsc{Rorario} sets down the page he has been holding.
\end{action}

\transition{CUT TO:}

% ------------------------------------------------------------
\scenenumber{16a — Animal Sequence: The Courtyard}
% ------------------------------------------------------------

\sceneheading{EXT. Study Courtyard — Afternoon}

\begin{action}
The courtyard visible from the study window. No human figures. The light at this hour produces a particular quality of stillness in open stone spaces: the heat absorbed by the surface during the earlier part of the day is now being released at a rate that makes the air immediately above the ground slightly warmer than the air at head height, a temperature gradient that is invisible but that insects and small animals navigate with precision, their flight paths and ground-level movements calibrated to its contours.

A \textsc{Sparrow} lands on the courtyard's near edge and moves across the stone in the rapid, interrupted progress characteristic of its species: advance, pause, advance, a tight turn, pause, advance --- the pauses being the intervals of intense sensory assessment and the advances being the committed consequences of that assessment, the whole sequence constituting a form of navigation through a space that the bird reads at a resolution of detail unavailable to the human eye. It locates something in a crack between two stones, extracts it with a single precise insertion of the beak, and departs at a diagonal angle that takes it out of the courtyard at a point determined by whatever factors govern the bird's sense of the optimal exit vector. The entire episode has taken perhaps thirty seconds.

Two \textsc{Pigeons} occupy the courtyard's center, moving through their characteristic slow-paced investigation of the ground surface, their heads bobbing with each step in the visual stabilization mechanism that their locomotion requires, the bobbing having the paradoxical quality of a compensatory stillness within movement, the eye held fixed relative to the world while the body advances beneath it. They find nothing of particular interest and continue their circuit without urgency.

The study window is visible from this angle as a dark rectangle in the building's face, the interior too dim at this hour to reveal what it contains. From within the study the courtyard is visible and animated. From the courtyard the study is opaque.
\end{action}

\transition{CUT TO:}

% ------------------------------------------------------------
\scenenumber{16d — The Staged Argument}
% ------------------------------------------------------------

\sceneheading{INT. Temporary Stage — Evening}

\begin{action}
A raised platform within a large interior, its boundaries demarcated by a slight elevation and a change in flooring material rather than by walls, the performance space defined by convention rather than enclosure. Seating has been arranged in a semi-circle, though not all positions are occupied. Three figures stand upon the stage in a triangular placement, each oriented slightly toward the center without directly facing the others --- a geometry that allows each to address both the other speakers and the audience without committing to either as primary addressee. A \textsc{Moderator} stands at the periphery holding a document that is not consulted.
\end{action}

\says{First Speaker}{The act was completed.}

\says{Second Speaker}{The act was prevented.}

\begin{action}
A brief interval. The Third Speaker does not immediately respond, the delay having the character of consideration rather than hesitation.
\end{action}

\says{Third Speaker}{The act was completed in its prevention.}

\begin{action}
The audience shifts slightly --- not in reaction but in accommodation, as though making room for an additional proposition without displacing the prior two, which remain in play as positions rather than as arguments that have been answered. The First Speaker steps forward; the Second replies without advancing; the Third turns slightly, addressing neither directly. Each elaborates their position through increasingly precise qualifications, none of which exclude the others entirely, though each attempts to do so through the progressive narrowing of the terms through which the act in question is characterized. At a certain point the First Speaker adopts a phrase introduced by the Second; the Second incorporates a term first used by the Third; the Third repeats a formulation from the First. The positions begin to share language without converging in meaning, the same words now performing different functions in different argumentative structures.

The triangular arrangement persists. No speaker occupies the center. The center remains empty --- not as absence but as a position that no one has taken, whose vacancy is a structural feature of the configuration rather than an oversight. The Moderator does not adjudicate. The argument does not conclude. It continues as a stable configuration of non-resolution, the stability consisting not in any agreement among the positions but in the mutual inability of each to displace the others.
\end{action}

\transition{CUT TO:}

% ------------------------------------------------------------
\scenenumber{17e — Horses in Rain}
% ------------------------------------------------------------

\sceneheading{EXT. Roadway — Rain}

\begin{action}
Rain falls with sufficient intensity to alter the surface of the road, converting previously firm ground into a variable substrate whose resistance to pressure shifts with each step, the alteration not uniform across the road's width but concentrated in areas where the drainage is insufficient and the water accumulates before dispersing laterally. Several horses are being led along the roadway, their movement continuous but uneven, the unevenness not produced by any failure of the animals' ordinary gait but by the road's continuous variation in the conditions it presents to each footfall.

A hoof meets the ground. It sinks slightly --- the ground yielding under the pressure to a depth determined by the local water content of the soil at that point, which differs from adjacent points in ways not predictable from the surface appearance --- and then stabilizes as the soil's resistance below the initial yielding layer provides sufficient support. The next step does not reproduce the prior condition; it encounters a different point on the road's variable surface, with different properties, requiring a different resolution.
\end{action}

\begin{action}
The handlers adjust their pace in response to the road's condition; the horses adjust independently, each responding to its own proprioceptive information about the ground rather than to any signal from the handler or the other horses. No single rhythm governs the group. A cart follows behind, its wheels tracing grooves that fill with water immediately after the wheels have passed, the groove present for a brief interval and then absorbed into the road's general surface condition, the track not remaining as a guide for subsequent passage.

One horse slips --- not fully, the slip arrested in its early phase by the redistribution of weight across the remaining limbs, the correction executed within the time available before the slip could develop into a fall. The correction introduces a minor deviation into the forward trajectory as the correcting movement carries the horse slightly to one side. The deviation propagates through the group as a series of local adjustments --- the following horses responding to the altered spacing, the handlers adjusting their positions, the cart shifting its line slightly. The group stabilizes into a new configuration that differs from the prior one in several of its parameters and is functionally adequate in all of them.

Rain accumulates along the road's edges in channels that redirect it unpredictably, each channel forming and redirecting itself as the accumulation exceeds the ground's absorption capacity at different points. The horses step through these without anticipation, each encounter resolved as it occurs rather than navigated in advance. No path is optimal under these conditions. No path is maintained across more than a few steps. Movement persists through continuous correction applied to a substrate that continuous correction cannot stabilize, the progression requiring not stability but sufficiency, which is a different and lower criterion that the conditions continuously meet.
\end{action}

\transition{CUT TO:}

% ------------------------------------------------------------
\scenenumber{10b2 — The Gate}
% ------------------------------------------------------------

\sceneheading{EXT. Courtyard Gate — Late Afternoon}

\begin{action}
A narrow gate in a stone wall, its width sufficient for a single cart to pass but not for two to pass simultaneously. A cart approaches from one side drawn by a horse; from the other, a second cart approaches, slightly larger, drawn by two oxen. The passage cannot accommodate both. The drivers perceive each other at a distance sufficient to allow adjustment. They slow.

One driver gestures. The other returns the gesture. Neither proceeds. They begin to speak. The exchange is polite, its early phrases the courteous formulations that such situations conventionally generate, followed by something more deliberate as each proposes a sequence of movement that would allow both to pass. The proposals are reasonable: each involves one cart reversing a short distance to allow the other to clear the gate first. The proposals require only coordination, not concession.

They do not move. One driver dismounts, measures the gate's width with his hands, steps off the distance between the carts, returns, makes a further proposal. The other considers this, begins to suggest a modification, pauses as a new consideration presents itself. Both drivers are stopped. One begins to reposition his cart and stops when the other simultaneously begins to do the same; the simultaneous initiation produces a mutual arrest. A third cart has appeared at the far end of the road and is slowing.

A boy stands nearby holding a rope attached to a dog. The dog attends to the gate opening with the quality of attention that a clear passage generates in an animal with a destination. The boy's grip relaxes slightly, or the dog pulls with a force that slightly exceeds his resistance. The rope passes through his hand. The dog moves into the gate opening, adjusts its body once to the available width, and emerges on the other side. It continues. The drivers, neither of whom has looked toward the dog, continue their exchange.
\end{action}

\transition{CUT TO:}

% ------------------------------------------------------------
\scenenumber{16f — The Ledge}
% ------------------------------------------------------------

\sceneheading{EXT. Rock Ledge — Day}

\begin{action}
A narrow ledge of stone, its surface weathered into a slight concavity. Seeds have been scattered across it by some prior process --- wind, or the disruption of a prior visit, or the opening of a seed casing above the ledge --- their distribution irregular in a way that reflects no single organizing principle.

A bird lands, takes one seed from the ledge's near end, and departs. A squirrel approaches from below, reaches the surface, moves across it while gathering two seeds, departs in a direction that differs from the bird's while leaving one seed behind at the point where it paused longest. Another bird arrives on the ledge and does not take the remaining seed, attending instead to the stone surface itself, scratching at some feature of the rock not visible from the available angle. A third animal --- small, its species not clearly established within the duration of its appearance --- emerges from a gap in the stone near the ledge's far end, takes the remaining seed, and disappears. The first bird returns to the ledge's near end, pauses at the point from which it previously departed, then moves along the ledge's length to a position it has not previously occupied. No trajectory repeats exactly. Each action completes within its own terms.
\end{action}

\transition{CUT TO:}

% ------------------------------------------------------------
\scenenumber{16g — The Water Edge}
% ------------------------------------------------------------

\sceneheading{EXT. Shallow Stream — Evening}

\begin{action}
A section of shallow stream, its bed visible through water whose current produces surface movement without obscuring the bottom. A horse stands at the stream's edge and drinks, the drinking interrupted at intervals by the horse lifting its head for a period that appears to be its own calibration of how much it needs.

Upstream and slightly visible within the same frame, a dog steps into the water. The step disturbs the stream's surface, the disturbance propagating downstream as a subtle alteration of the current. The horse lifts its head, shifts its position by a small distance along the bank, and drinks again. Whether the shift is a response to the upstream disturbance or an independent adjustment of its standing position is not determinable from the available evidence.

A bird skims the water surface and takes something from it --- something small and at or just below the surface, the extraction so rapid that the object itself is not identified in the motion. It continues along the stream's course and disappears from the frame. The current carries a small piece of debris past all three animals and past the observation point. Each action is complete within its own terms. No interaction resolves into a stable relation between any pair of the three presences.
\end{action}

\transition{CUT TO:}

% ------------------------------------------------------------
\scenenumber{16h — The Branch Network}
% ------------------------------------------------------------

\sceneheading{EXT. Tree Canopy — Day}

\begin{action}
Several branches of a tree, their arrangement producing a network of intersecting paths at a scale that birds navigate by adjusting their flight to the available intervals between branches. Several birds are present, each on a different branch or moving between them, each attending to something on its branch or adjusting its position relative to an adjacent bird or moving to a new branch with the directedness of a bird moving toward something rather than simply changing position without reason.

One bird drops a fragment from its bill. Another takes it before it has fallen more than a short distance, the interception occurring within the same continuous motion that brought the second bird to that position. A third bird lands on the branch where the fragment had been. For perhaps three or four seconds, the positions and movements of the birds exhibit a spatial regularity that suggests a pattern --- something like a circuit, something like a sequence, something like the distributed action of a system that has a logic to it.

Then one bird departs in a direction that breaks the regularity. Another repeats a movement that leads to nothing, a branch-position that generates no further action. The apparent circuit does not close. The birds continue in the canopy, each attending to its own immediate situation, the suggestion of a pattern neither confirmed nor denied by what follows, simply no longer generating the appearance that briefly proposed it.
\end{action}

\transition{CUT TO:}

% ------------------------------------------------------------
\scenenumber{16i — The Mechanism}
% ------------------------------------------------------------

\sceneheading{EXT. Edge of Wood — Day}

\begin{action}
A sprung trap hangs slack at the edge of the wood, its mechanism having been triggered at some prior point, the bent branch now released. A strip of dried meat remains caught in the loop, positioned just above the ground.

A squirrel approaches the mechanism from the tree side, descending a trunk and moving along a lower branch that brings it within proximity of the loop. It tests the branch with its weight; the branch dips slightly; the meat shifts within the loop but does not fall. The squirrel withdraws along the branch and moves to a higher position in the tree.

A bird lands on a branch near the mechanism, watches the loop without approaching, does not interact with it. The squirrel moves again along a different branch, carrying something small. It releases the object; the object falls and strikes a branch above the trap; the branch trembles; the meat shifts again without falling. The squirrel descends partway and reaches toward the loop, then withdraws abruptly.

Wind passes through the edge of the wood. The branch moves. The meat slips, rotates once within the loop, and falls to the ground.

Neither animal moves immediately. After a period of stillness, the squirrel descends to the ground and takes the meat, moves a short distance, and stops. The bird drops from its branch to the ground near where the meat had been, pecks once at the surface, and returns to the branch. The squirrel eats. The bird remains. Nothing passes between them. The mechanism hangs empty.
\end{action}

\transition{CUT TO:}

% ------------------------------------------------------------
\scenenumber{17 — The Visitor}
% ------------------------------------------------------------

\sceneheading{INT. Study — Day}

\begin{action}
A \textsc{Woman} visits. She is a relative, or someone from an earlier period of \textsc{Rorario}'s life --- the film does not specify, because the specification is not available from the evidence that the scene presents, which is the evidence of two people in a room together, their shared history present in the ease of certain small gestures and absent from the verbal content of their exchange, which is conducted at the level of present circumstance rather than accumulated past. She has come from within the city; she does not carry the marks of travel.

The conversation is warm in its register, which is to say it occupies a different tonal range than any of the conversations that have preceded it in the film without constituting a different kind of event. She speaks. He responds. She is from the prior system in the sense that she inhabits without apparent difficulty the assumption that conversation is a medium for the transmission of content from one mind to another, that words mean what they mean in the context of a mutual attention that both parties can sustain, that the exchange of statements constitutes a form of contact between persons. She proceeds on these assumptions without checking them.

Her words reach \textsc{Rorario} and are processed, for a time, with the normal machinery of reception. He asks a question. She answers. The exchange has the natural rhythm of two people for whom conversation in this register is not effortful.
\end{action}

\begin{action}
At some point during her account of something --- a family matter, an event in the city, the condition of a mutual acquaintance --- her words begin to arrive at \textsc{Rorario} in a mode that is subtly different from their earlier arrival. The words continue to be spoken; their prosodic features continue to be received; the syntactic structures continue to be parsed at the level of form. But the semantic content --- the propositional substance, the meaning that the words are transmitting --- undergoes a reduction in the depth of its processing that produces, not incomprehension, but a kind of reception-without-elaboration: the words arrive and are registered without generating the inferential extensions and associated representations that would constitute understanding them fully. The sound of speech without the weight of speech. The formal envelope of communication without its content.

This is not marked by any change in \textsc{Rorario}'s behavioral presentation. He continues to maintain the postures of attentive reception. He continues to respond at the appropriate intervals. The responses are adequate to the surface structure of what she has said. She does not perceive a change.

The perceptual shift is not restored until after she has risen to leave, gathered her things, and moved through the doorway. In the interval immediately following her departure, the room reassembles its acoustic character without her voice in it, and \textsc{Rorario}'s reception of that character --- the ambient sounds of the building, the distant courtyard, the movement of air through the study's single window --- returns to the fuller depth of processing that the conversation had briefly suspended.

The film does not protect her from its formal logic. The reweighting of signal that has been building across State~II has now extended into the register of intimacy. Nothing in the intimacy of the exchange exempts it.
\end{action}

\transition{CUT TO:}

% ------------------------------------------------------------
\scenenumber{17a — Isaac's Ashes: A Fragment}
% ------------------------------------------------------------

\sceneheading{INT. Apothecary's Workshop — Day}

\begin{action}
The workshop is a space of accumulated material: dried plants in bundles suspended from low beams, vessels of varying size and opacity arranged on shelving whose organization follows a logic the practitioner knows but has not made externally legible, a work surface bearing the evidence of recent procedures --- residues at the rims of vessels, a fine pale powder at the center of the surface that has been partially displaced by whatever activity has most recently been conducted here but has not been fully cleared.

An \textsc{Apothecary} works at the surface, engaged with the reduction of a substance by heat and grinding. The procedure is of the kind that requires sustained attention of a particular quality: not intense but continuous, the monitoring of a process rather than the execution of a discrete act, the practitioner present to a transformation that is occurring at its own pace and cannot be rushed without being altered.

\textsc{Rorario} has come with a small commission whose nature is not specified by the scene. He waits. The Apothecary works. From an adjacent room, the voice of an \textsc{Assistant} can be heard reading aloud from a text --- a devotional or instructional passage whose content arrives in the workshop as fragments, the wall between the rooms absorbing the middle of each sentence while allowing its beginning and end to pass through:
\end{action}

\begin{action}
\textit{\ldots{}the ashes upon the mountain, which were gathered\ldots}

\textit{\ldots{}and God, who restores\ldots}

\textit{\ldots{}as from the heap of what was consumed, the renewed\ldots}
\end{action}

\begin{action}
The fragments arrive in this form: not as a continuous text, not as a doctrine, but as pieces of a statement that is being made in full in the adjacent room and is arriving here in a state of partial transmission. Their meaning is neither clear nor unclear; they carry the impression of a subject without specifying it, the way that overheard speech at a distance carries the formal features of significance without making its content available.

The Apothecary does not look up from the work surface. The pale powder at the center of the surface is the residue of the heating process he has been conducting. He gathers it carefully with a small instrument and transfers it to a vessel, which he seals. The gathering is complete. He turns to \textsc{Rorario}.

The reading in the adjacent room has continued and has moved beyond the passage whose fragments were audible here. What is audible now is simply the cadence of reading, the words no longer individually distinguishable at this distance.
\end{action}

\transition{CUT TO:}

% ------------------------------------------------------------
\scenenumber{17b — Translation Chamber}
% ------------------------------------------------------------

\sceneheading{INT. Translation Chamber — Late Afternoon}

\begin{action}
The chamber is narrow, its walls lined with textiles whose principal function is acoustic rather than decorative --- the dampening of reflected sound so that the spoken word, in this space of double articulation, arrives at the receiving desk without the distorting reverberance that stone surfaces introduce. Two desks face one another across a confined interval, each occupied by someone whose task is the relay of meaning across a linguistic boundary that is itself of indeterminate width. A document lies between them, oriented such that it may be read from either side without being physically moved --- a positioning that performs neutrality while being inevitably more legible from one angle than the other, the original language being, by the nature of the situation, always the language of the document rather than the language toward which it is being translated.

A \textsc{Speaker} reads from the original text. A \textsc{Translator} renders it into another language, his production following the Speaker's by an interval that varies slightly with the syntactic complexity of each clause, the variation registering the different cognitive load imposed by clauses that map cleanly across the linguistic boundary and clauses that resist clean mapping and require the Translator to make small decisions about which of the available target-language structures best approximates the source's intent.
\end{action}

\says{Speaker}{The boundary shall be observed where it has been placed.}

\says{Translator}{The boundary shall be observed where it is recognized.}

\begin{action}
The substitution of \textit{recognized} for \textit{placed} introduces a shift in the boundary's ontological condition: from a physical fact about the world to an epistemic condition of its observers. Whether the boundary exists where it was placed, or where it is recognized to be, are not the same question, and the difference between them is not a matter of phrasing but of the entire framework within which the boundary's authority is grounded. The Translator has made a decision about which framework is appropriate for the target language's legal and administrative register, and the decision is locally defensible. No objection is raised.
\end{action}

\says{Speaker}{Recognition shall follow from its establishment.}

\says{Translator}{Recognition shall establish it.}

\begin{action}
The inversion is syntactically minimal. Its consequence is not. The original formulation preserves the primacy of physical establishment as the ground of recognition; the inverted formulation makes recognition itself constitutive, which is a different theory of how boundaries come to exist. The Translator has not introduced the inversion maliciously or carelessly; he has produced the formulation that his target language's administrative idiom most naturally generates when the source statement is translated into it, and the natural generation of the target idiom happens to produce a different theory. This is a structural feature of the translation relationship, not a personal failure of the Translator.

A \textsc{Clerk} at the side table copies the translated version without consulting the original. He writes with steady pace. The two texts diverge incrementally, each step preserving local coherence while the cumulative drift exceeds the tolerance of correspondence. No correction is made. None is requested. The copied document is sanded, folded, and sealed. It is dispatched. The original remains. The two texts now coexist, no longer equivalent, each capable of further propagation into the network of subsequent documents that will cite them as authoritative sources for what the boundary is.
\end{action}

\transition{CUT TO:}

% ------------------------------------------------------------
\scenenumber{17c — The Cartographic Room}
% ------------------------------------------------------------

\sceneheading{INT. Cartographic Room — Day}

\begin{action}
The room contains a large table upon which a map is spread and partially weighted at its corners by objects whose function is purely ballistic --- small stones, a book, a knife whose blade has been used for this purpose so frequently that its edge bears the marks of the stone's surface. The paper has been repaired in several places, seams visible where sections have been replaced or reattached when the original material proved insufficient for the territory it was required to represent, the repair itself introducing slight misalignments at each seam as the new paper and the old are brought into approximate rather than exact correspondence.

Lines traverse the surface in several registers: survey lines in one ink, administrative boundaries in another, and a third set of lines in a third ink whose age is difficult to determine and whose relation to current conditions is consequently uncertain. Some lines are continuous; others are interrupted and resumed at slight offsets across the repair seams, the interruption and resumption together producing the visual effect of continuity while introducing, at each seam, a small angular deviation whose cumulative effect across the map's extent is not negligible.

\textsc{Rorario} stands over it. A \textsc{Cartographer} indicates a region.
\end{action}

\says{Cartographer}{This corresponds to the current boundary.}

\begin{action}
He traces a line with his finger, a gesture that is both clarifying and slightly distorting, the finger's width covering the fine detail of the line's actual course while the gesture's confidence implies a precision the line does not quite possess. \textsc{Rorario} follows the indicated course with his gaze.
\end{action}

\says{Cartographer}{This corresponds to the revised boundary.}

\begin{action}
A second line, nearly parallel to the first, intersects it at two points and diverges between those intersections and beyond them, the interval between the lines not constant but varying with the topographic features that each line tracks according to different principles of boundary determination. The map does not privilege either line. Both are inscribed with equal weight of ink, equal authority of draftsmanship, equal apparent confidence about the territory they represent.

A section of the map in the area between the two lines is missing, replaced by a patch of different paper. The lines continue across it, but their alignment shifts slightly at each seam of the repair, the accumulated deviation at the patch's far edge placing each line in a position slightly different from the position it would have occupied if the map had been drawn on a single continuous sheet. The Cartographer smooths the surface with his palm as though the gesture might settle the seams into better correspondence. The lines do not reconcile.

\textsc{Rorario} does not request clarification. He places a small marker on the map. It falls between the two lines, in the interval that the map represents but does not resolve. The Cartographer nods. The placement is recorded in a separate document, which will subsequently be consulted by parties who do not have access to the map and who will therefore receive the recorded placement as a precise coordinate rather than as a position within an unresolved interval.
\end{action}

\transition{CUT TO:}

% ------------------------------------------------------------
\scenenumber{17d — The Corridor Plant}
% ------------------------------------------------------------

\sceneheading{INT. Administrative Corridor — Midday}

\begin{action}
The corridor connects several offices within the same administrative complex, its function purely transitional --- a space through which one passes in order to arrive at spaces designated for other activities. It has no designated function of its own and consequently receives no designed attention: its surfaces are maintained to the standard of shared transitional spaces, which is to say they are maintained, but not to the standard that any individual office would maintain a space it occupies with sustained presence. In this condition of institutional inattention, a plant has grown from a fissure near the base of the wall on the corridor's north side.

The plant is of a succulent type, its stems thickened with stored water, its leaves broader than its anchoring point in the fissure's accumulated debris would seem capable of supporting. It has been growing for some time; this is evident from its extent, which reaches into the passable width of the corridor far enough to require minor adjustments in path from anyone who passes along the wall's side. The plant has not been placed here. It has not been authorized. It has not been prohibited. It occupies the corridor as an outcome of its own biological processes, which have followed the available resources --- the slight moisture that seeps through the fissure, the light from a window at the corridor's end --- without reference to the administrative character of the space they have colonized.
\end{action}

\begin{action}
A \textsc{Clerk} carrying documents approaches along the wall side. He shifts slightly to avoid the leaves, the adjustment so minimal and so habituated that it does not register as a decision, only as a local modification of his standard path. A page brushes against the leaves at the outermost extent of his path-shift. A faint green transfer marks the margin of the page --- a pigment deposit from the leaf surface, concentrated at the point of contact, fading toward the page's text. The Clerk does not notice the transfer. He continues. The document bearing the transfer will be read by its recipient in a room where no plant is present, and the green mark at its margin will be attributed to whatever local explanation presents itself, or will not be noticed, or will be noticed and found unremarkable.

Another figure passes. He pauses briefly, touches one of the leaves with an index finger in the absent gesture of someone registering the presence of something that does not require a response, and proceeds. The contact leaves no visible change in the plant. It remains in its current configuration, neither enlarged nor diminished by the touching.

\textsc{Rorario} passes through the corridor. He does not alter his path. A leaf brushes his sleeve at the same point of the corridor where it brushes every sleeve that passes on the wall side, the plant's position having become, through repeated contact, a known feature of this particular passage. The contact is registered without being processed. He continues. The plant's presence in the corridor has not been formally acknowledged at any administrative level. It has been incorporated into the corridor's operational character as thoroughly as if it had been designed into it.
\end{action}

\transition{CUT TO:}

% ------------------------------------------------------------
\scenenumber{19a — Contracting Letters}
% ------------------------------------------------------------

\sceneheading{INT. Writing Chamber — Evening}

\begin{action}
The chamber contains several writing surfaces, each occupied by someone engaged in the production of correspondence whose destinations are various and whose internal relationships to one another are partially legible from the pattern of documents being consulted as reference material. The light is insufficient for sustained clarity at this hour; candles have been brought closer to the individual writing surfaces, producing an arrangement in which each desk is surrounded by its own small zone of adequate illumination separated from adjacent zones by intervals of relative dimness, the chamber now a distributed system of local light rather than a unified field.

A \textsc{Clerk} writes at the central desk. The letter he is producing has been underway for some time; its opening has been drafted, crossed out, and redrafted, though the current version retains elements of the original formulation that the crossing-out was intended to eliminate, the words surviving in the new draft through the mechanism by which formulations that have been used for some purpose tend to recur in subsequent formulations of the same content.
\end{action}

\begin{action}
He rereads a line. The phrase \textit{as previously established} appears in the draft at a point where the letter must refer to an arrangement that was itself the product of an earlier correspondence now in circulation elsewhere in the administrative network. The phrase performs its function of establishing continuity with prior material without specifying which prior material, because the writer knows which prior material and assumes the reader will also know, an assumption that has been correct in prior instances and may or may not be correct in the present instance, since the relevant prior material exists in several versions of which this writer is aware of some and the intended reader of others.

He continues. The sentence extends through additional clauses, each of which is locally clear and each of which increases the sentence's internal cross-referential density to a point at which the clause currently being written refers to a condition established by a preceding clause that itself refers to a prior document, the chain of reference now extending through two links whose stability cannot be verified within the writing act itself. He pauses. He rereads the sentence from its beginning. He reaches the disputed clause. He continues past it, adding a further clause that partially qualifies the preceding one without resolving the ambiguity it contains.

At an adjacent desk, another \textsc{Clerk} writes more quickly, his sentences shorter, their internal cross-referential structure reduced to a minimum that allows the correspondence to function without requiring the full apparatus of verification that longer sentences would imply. His letters omit qualifiers that the letters of the first Clerk include. The omission simplifies without falsifying; it also removes from the record the cautious hedging that careful administrative writing uses to acknowledge the limits of its own certainty. His letters are easier to act upon than the first Clerk's. What they are easier to act upon toward is less precisely specified.
\end{action}

\begin{action}
A \textsc{Messenger} waits at the chamber's margin for the letters to be completed. His waiting does not produce any acceleration in the completion of the letters; he is a feature of the chamber's operational environment rather than a pressure applied to it. The first Clerk sands the page, examines it in the candlelight, and before folding it adds a further clause in the margin --- a qualification of the document's primary instruction that is necessary for precision but that, by being marginal rather than integrated, introduces a new ambiguity about the instruction's relationship to its qualification, the marginal clause being legible as either a condition on the primary instruction or a supplement to it, depending on the interpretive conventions of the reader.

He folds the letter. He hesitates. He unfolds it and reads the marginal clause again. He folds it. He seals it. He hands it to the Messenger, who receives it along with the shorter letter from the adjacent desk. The Messenger departs with both. Their relation to one another --- whether they are companion communications intended to be read together or independent documents whose simultaneous dispatch is incidental --- is not specified by anything in their external presentation. They will arrive together. They will be read separately. Their combined effect on the situation they address will be neither the effect of either alone nor the sum of both, but something produced by the interaction of their respective formulations in the mind of whoever must act upon them.
\end{action}

\transition{CUT TO:}

% ------------------------------------------------------------
\scenenumber{15c — Tavern: Goliath Rotating Roles}
% ------------------------------------------------------------

\sceneheading{INT. Tavern — Night}

\begin{action}
The tavern interior is thick with suspended particulates --- smoke, breath, the residue of prior occupation not yet dissipated into the inadequate ventilation of a room designed for other purposes than air quality. A makeshift platform has been assembled from two tables drawn together, their surfaces uneven, the unevenness corrected by wedges of folded cloth placed at the legs that require them, a practical adjustment that gives the platform the character of something improvised within constraints rather than designed from the outset.

A small assembly gathers, not formally seated but cohering around the platform through the minor adjustments of stance and orientation that constitutes an audience in a space not designed to produce one. They have collected from the room's ambient population through the mechanism by which performances in public spaces generate their audiences: proximity, then attention, then the slight social pressure of being already among those who are watching.
\end{action}

\begin{action}
Two men stand upon the platform. One is markedly larger than the other, the difference in scale sufficient to render the dramatic relation between them legible without supplementary identification. The other, not.
\end{action}

\says{Large Man}{I am Goliath.}

\says{Smaller Man}{I am David.}

\begin{action}
A third man, already present in the assembly but not initially central to its attention, steps forward and places a sling into the hand of the Smaller Man, the instrument arriving without ceremony, as though its availability at this moment were unremarkable. The scene proceeds along its established course: the Smaller Man casts, the Large Man falls with a degree of deliberation that slightly exceeds what impact would require, as though the fall were being performed to the specification of its expected form rather than produced by its cause.

A brief interval. The fallen man rises. He crosses to the other side of the platform and exchanges positions with a member of the assembly who has moved forward to receive him.
\end{action}

\says{New Large Man}{I am Goliath.}

\says{Former Large Man}{I am David.}

\begin{action}
No announcement accompanies the substitution. The assembly adjusts its understanding of the scene's current personnel without visible confusion, the reassignment of roles accepted as a feature of the form rather than a deviation from it, as though the identity of Goliath and the identity of David were understood to be structural positions within the event rather than properties of the specific bodies that occupy them at any given iteration.

The sling is passed again. The sequence repeats. The relation between role and body continues to de-stabilize by increments, each exchange making it more apparent that the role adheres to a position within the event's structure rather than to any particular person, and that the event's structure is itself the stable element that the rotating personnel demonstrate rather than undermine.
\end{action}

\begin{action}
At one iteration, two men step forward simultaneously.
\end{action}

\says{Man One}{I am David.}

\says{Man Two}{I am David.}

\begin{action}
A brief hesitation in the assembly's collective attention, as though the duplication had introduced a condition the form's usual rules do not address. A voice from the periphery:
\end{action}

\says{Voice (O.S.)}{Then one must be Goliath.}

\begin{action}
Neither man yields his designation. The scene continues regardless, a fall occurring without the duplication having been resolved, the assembly registering the outcome without requiring that the premise be settled before the consequence can be accepted. The platform remains. The roles continue to circulate. The narrative persists in the absence of a stable assignment of its principal positions to specific persons. This does not appear to impair its operation as an event.
\end{action}

\transition{CUT TO:}

% ------------------------------------------------------------
\scenenumber{15d — Children's Game: Isaac Reset}
% ------------------------------------------------------------

\sceneheading{EXT. Courtyard — Afternoon}

\begin{action}
A group of children has assembled within a bounded but unenclosed space, its definition produced less by any architectural feature than by the habitual convergence of activity in this particular area, which has acquired through use the character of a designated space without having been designated. The ground bears marks of prior games --- lines partially effaced by subsequent use, circles whose boundaries intersect without any resolution of which circle has precedence having been established or required.

One child lies supine upon a raised surface composed of stacked stones, the surface elevated perhaps thirty centimeters above the ground, sufficient to distinguish it from the ground without constituting a significant height. Another child stands beside him, holding a stick whose function in the game's symbolic economy is indicated by the context rather than by anything inherent in the object.
\end{action}

\says{Standing Child}{You must not move.}

\says{Prone Child}{I will not move.}

\begin{action}
The Standing Child raises the stick. No evident force is applied or implied; the gesture is symbolic in the strict sense, its meaning residing in its position within the sequence rather than in any physical consequence it produces or threatens to produce. The Prone Child closes his eyes, the gesture of closing performing the required condition of non-witness to what is about to occur.

A third child intervenes, placing a hand upon the arm of the Standing Child at a point in the sequence that has apparently been established as the correct point for the intervention.
\end{action}

\says{Third Child}{Not yet.}

\begin{action}
The sequence halts at the intervention point. A discussion ensues among the assembled children, not about whether the act should occur --- that question is not in dispute --- but about whether the current moment constitutes the correct timing for the interruption, the rules of the game apparently specifying a particular temporal position for the intervention that is itself defined by some condition whose satisfaction the children are now adjudicating. The discussion is brief and procedural. After a moment the Third Child withdraws his hand.
\end{action}

\says{Third Child}{Now.}

\begin{action}
The stick is lowered. The Prone Child remains still through the duration of the gesture's completion. A pause follows, during which the event that has occurred is allowed to register as having occurred.
\end{action}

\says{Another Child}{He is dead.}

\begin{action}
Silence of the appropriate duration. Then:
\end{action}

\says{Third Child}{Again.}

\begin{action}
The Prone Child opens his eyes, sits up, descends from the stone surface, and exchanges places with another child without ceremony, the exchange conducted with the matter-of-fact efficiency of a procedure that has been repeated sufficient times to have been fully normalized. The structure repeats with the new Prone Child in position. No account is kept of prior iterations; each instance is treated as complete in itself, the preceding iterations having no bearing on the current one's validity.

At one iteration, the interruption fails to occur at the established moment --- the Standing Child lowers the stick before the Third Child's hand arrives. A brief uncertainty arises among the assembly about the status of this iteration: whether the failure of the interruption to occur at the correct moment constitutes a completed act or an incomplete one, whether the child who was the subject of the uncompleted interruption has the same status as those whose iterations proceeded correctly. The question is raised but not resolved before the child rises regardless, the game's continuity asserting itself as more pressing than the adjudication of the exceptional case. The game continues. The boundary of the courtyard remains wherever the repetition has established it.
\end{action}

\transition{CUT TO:}

% ------------------------------------------------------------
\scenenumber{15e — Puppet Theater: Cain and Abel}
% ------------------------------------------------------------

\sceneheading{INT. Puppet Theater — Evening}

\begin{action}
A small enclosed structure houses a raised stage upon which figures of wood and cloth are arranged, their construction simple enough that the convention of their representing human figures is established by their relative scale and their placement within the scene's action rather than by any realistic rendering of the human form. A curtain separates the operators' space from the performance space, though its lower edge admits glimpses of the mechanisms behind, the machinery of the performance partially visible at the point where the convention is thinnest. A \textsc{Narrator} speaks from behind the curtain. The puppets move.
\end{action}

\says{Narrator}{The elder brother brought the offering that was not accepted.}

\begin{action}
On the stage, the puppet corresponding to the elder places an object upon a small altar constructed of painted wood. The object is received by the altar --- it remains upon it, upright, stable. It is not rejected. The Narrator's account and the puppet's action are in this respect discrepant: the narration describes an outcome that the performance does not enact.
\end{action}

\says{Narrator}{The younger brought that which was favored.}

\begin{action}
The second puppet presents its offering. It falls from the altar and separates from itself on the stage floor, the construction of the puppet prop having been insufficient to withstand the force of its placement or the angle of its contact with the altar's surface. A murmur passes through the audience, not of alarm but of the quality of attention that unexpected deviations from expectation produce in those who are following both a narration and its enactment simultaneously. No interruption occurs. The performance continues.
\end{action}

\says{Narrator}{And the elder was marked.}

\begin{action}
A mark is applied to the stage itself rather than to the elder puppet --- a darkened area produced by some property of the staging materials, spreading beneath both figures without discrimination between them. The puppets continue to move within it. A string tangles; one puppet jerks irregularly, the jerk introducing an action not present in the narration, which continues without adjustment to account for it. The narration describes an order; the enacted sequence does not correspond to that order. Both are maintained simultaneously. The audience attends to both without the performance providing guidance about which constitutes the authoritative account of what is occurring.
\end{action}

\says{Narrator}{The order is preserved.}

\begin{action}
The curtain lowers. The discrepancy between the narrated and the enacted is not resolved at the performance's conclusion; it is completed, which is to say the performance ends while the discrepancy remains, the ending being a formal property of the event rather than a consequence of the discrepancy's resolution.
\end{action}

\transition{CUT TO:}

\sceneheading{EXT. Open Field — Late Afternoon}

\begin{action}
A large aggregation of birds occupies the air above an open field at that hour of the afternoon when the light is horizontal enough to catch the undersides of wings, producing brief flashes of brightness as the formation turns. Their motion is continuous --- it does not begin within the scene's available observation and will not end within it; the scene arrives upon a process already underway and departs from it still underway, the process having no natural terminus within the scale of observation that the scene provides.

No single bird appears to direct the movement. The formation's changes of direction, its contractions and expansions, its rotations around axes that are themselves moving, are produced through the aggregate of local adjustments, each bird responding to the movements of its proximal neighbors rather than to any signal from a central point. The coherence of the whole is a property that emerges at the scale of the whole, not a property that exists at the scale of any of its components, which are simply birds adjusting their flight paths in response to the birds immediately adjacent to them.
\end{action}

\begin{action}
The formation contracts. A compression moves through it in a wave whose propagation speed exceeds the speed at which any individual bird has moved, the wave being a property of the collective rather than of any constituent. The formation expands. It rotates along an axis that is not fixed in advance but is produced through the simultaneous adjustments of thousands of birds, each of which is making a local decision that contributes to a global direction that no individual bird has decided upon. The sound present is not a discrete signal but a distributed modulation arising from the aggregate of wingbeats, a sound that has frequency and amplitude but not the structure of communication, accompanying the motion without directing it.

At the field's edge, two figures stand and observe. They do not speak. The birds descend briefly toward the field's surface and rise again, maintaining the coherence of the formation across the interruption of the descent without any apparent coordinating mechanism, each bird's rise producing the conditions for the rise of adjacent birds in a cascade that restores the formation's altitude without requiring any bird to have planned the restoration. No record is made of the formation's prior states. The formation does not require its prior states to be remembered in order to continue. It is sufficient that it continues.
\end{action}

\begin{action}
The pattern persists beyond the duration for which any observer can retain its prior configurations in memory, the present state of the formation no longer being comparable to any retained prior state because the prior states have been replaced by subsequent ones before comparison can occur. This is not a failure of the observer's memory. It is a feature of the event: it has a scale of temporal complexity that exceeds the retention capacity of any single observer, and it is nonetheless entirely coherent. The two figures at the field's edge continue to observe. They have nothing to say about what they are watching that would add to it.
\end{action}

\transition{CUT TO:}

% ------------------------------------------------------------
\scenenumber{18a — Boundary Revisited}
% ------------------------------------------------------------

\sceneheading{EXT. Boundary Zone — Dusk}

\begin{action}
The boundary site from an earlier scene reappears under the altered conditions of a later hour and a later season, both differences registering primarily in the quality of the light and the state of the vegetation, which has grown into the interval between the two original markers in ways that neither survey team anticipated and neither administrative record reflects. The markers remain, but their relation to one another has been modified --- not through any direct intervention upon the markers themselves, but through the accumulation of minor interventions in their vicinity, each of which responded to the conditions it found rather than to the conditions that the original placement had intended to produce.

Additional stakes have been introduced by parties whose identities are not established by anything visible in the scene. Some align with the trajectories established by the original markers. Others do not, having been placed in response to features of the ground, access requirements, or local practical necessities that the original survey did not account for. The field has become denser with markers than it was, and the density has not clarified the boundary so much as it has multiplied the number of positions in relation to which the boundary's location must be understood.
\end{action}

\begin{action}
\textsc{Rorario} stands at a point between two lines that now diverge more noticeably than when he last stood here, the divergence having been produced not by any sudden change but by the accumulation of the small adjustments that each party has made to its own line in response to the instructions that have reached them through their respective administrative channels, the adjustments individually minor and collectively significant. A \textsc{Surveyor} consults a document whose relation to the current state of the field is one of partial correspondence, the document having been produced from an earlier state and having been authoritative in relation to that state without having been updated to reflect subsequent developments.
\end{action}

\says{Surveyor}{This corresponds to the established boundary.}

\begin{action}
He indicates one line. Grass grows differently within the interval between the two lines, the difference in growth reflecting differences in the use to which the two sides of each line have been put, which in turn reflects the understanding of each adjacent party about which line constitutes the operative boundary, which understanding varies between parties in the manner already established by the film's prior treatment of the boundary's administrative history.

A third set of markers, more recently placed than either original set and of a different material, intersects both lines at irregular intervals without aligning with either. Their introduction has produced a further set of implicit boundaries --- the lines connecting one third-party marker to the next --- that now coexist with the two primary lines without being formally related to either.
\end{action}

\says{Surveyor}{This corresponds to the recognized boundary.}

\begin{action}
He indicates the other. \textsc{Rorario} steps forward, placing his foot within the interval between the two lines, his position within that interval not exactly at its center but at a point whose selection appears neither deliberate nor random --- a position produced by the available footing, the angle of his approach, the inclination of the ground. A \textsc{Clerk} records the position in a document, the act of recording requiring him to select among possible representations of the position's relationship to the two lines, a selection he makes with brief hesitation and executes with the confidence that documentation requires to appear authoritative.

In the distance, a group makes a slight adjustment to one of the third-party markers. The adjustment is small. Its propagation through the network of alignments that subsequent surveyors will establish in relation to it is not small. The boundary persists in the field as a condition --- a distributed feature of the landscape's administrative interpretation rather than a property of the landscape itself, present in the stakes and the documents and the divergent grass and the different uses of adjacent parcels, nowhere present as a single legible line.
\end{action}

\transition{CUT TO:}

% ------------------------------------------------------------
\scenenumber{19b — Record Room: Reference Collapse}
% ------------------------------------------------------------

\sceneheading{INT. Record Room — Night}

\begin{action}
The room is densely occupied by shelves containing documents of varying age and condition. The arrangement is systematic in principle, though successive additions have introduced local irregularities whose correction would require a reorganization that the room's ongoing operational demands do not permit. A \textsc{Clerk} retrieves a document from its designated location and compares it with another already open on the central surface. He reads from the retrieved document.
\end{action}

\says{Clerk}{As clarified in the prior instruction ---}

\begin{action}
He searches for the referenced instruction. It is located in a different section of the shelving. The phrasing differs from what the referencing document implies it will say: not contradictory, but operating at a different level of specificity, its clarification addressing aspects of the matter that the referencing document's use of the phrase \textit{prior instruction} does not specify it addresses. A third document is introduced, one that references both of the prior two. Its terminology aligns with neither, having been produced at a later date by a different office using a different administrative vocabulary that was not yet in use when the first two documents were written.
\end{action}

\says{Clerk}{As amended in the subsequent clarification ---}

\begin{action}
The Clerk attempts to establish a sequence --- original, clarification, amendment --- that would allow the three documents to be read as a progressive development of a single line of instruction. The sequence cannot be stabilized without excluding one document's claim to its designated position within it, since the chronological order and the logical order of the three documents are not identical, and neither order, when followed consistently, produces a fully coherent reading of all three.

He writes a summary incorporating elements from all three documents. The summary introduces a new phrasing not present in any of its sources.
\end{action}

\says{Clerk}{As previously coordinated ---}

\begin{action}
The phrase does not appear in any prior document. It is accepted as functional --- it performs the work of designating a prior state without specifying which prior state, a generality that allows it to refer to the composite of the three documents without requiring that their relation be resolved. Another Clerk consults the summary and copies it, the copy omitting one clause whose relationship to the rest the copying Clerk has assessed as redundant and whose omission reduces local ambiguity while increasing the global ambiguity of the summary's relationship to its sources. A Messenger waits for the documents. They are not yet aligned in the sense of being mutually consistent, but they are aligned in the sense of being ready for dispatch, which is the operative criterion. They are released. The references now circulate in their altered form, each iteration preserving operability while the relationship between them continues to drift from whatever relationship they had when they were first produced.
\end{action}

\transition{CUT TO:}

\sceneheading{INT. Study — Late Afternoon}

\begin{action}
The manuscript is present on the table in a configuration that differs from all its prior configurations: the pages are gathered into a single stack whose arrangement is stable, the edges of the pages approximately aligned, the overall object having the character of something that has arrived at a state that is not so much finished as no longer in active modification. The arrival at this state has not been marked by any completing act --- no final entry, no last page written with the awareness that it was the last, no colophon, no date. The manuscript has achieved its current state through the cessation of the process that had been producing it, which is a different thing from completion.

\textsc{Rorario} is at the table. He has set down the pen. Not just placed it aside as in the motion of someone who will return to it --- set it down in the way that signals a different relationship to the object than the one that had been operative during use. The pen is at the table's edge, not in the working area. The inkwell is closed. These are small but legible indicators of a transition in the room's operational state.
\end{action}

\begin{action}
The afternoon light enters from the window at the angle appropriate to this hour, lower and more horizontal than at midday, producing a different distribution of illumination across the table's surface than has prevailed during the writing sessions that have constituted this room's primary function in the film. The manuscript's gathered pages are illuminated across approximately two thirds of their surface; the remaining third falls into the mild shadow produced by the stack's own thickness in relation to the light's angle.

Outside, in the courtyard, two \textsc{Dogs} are present. One has been there for some time --- it is in the posture of a dog that has found a location and committed to it, its body settled into the stone's warmth with the specific quality of repose that follows a decision about where to be rather than the temporary stillness that precedes a change of location. It settles immediately: a single motion of lowering that does not require the preparatory circling that some animals perform before lying down, the commitment executed without negotiation.

The second dog arrives from the courtyard's far entrance. It moves across the courtyard with purpose, but the purpose, on arrival, proves not to be any specific destination --- it circles once, a medium-radius orbit of the general area where the first dog lies, neither spiraling inward nor maintaining a stable orbit but executing a single circumscription and then lowering itself to the ground at a point approximately one body-length from the first dog. The settling is complete within the same cut that contains the circling. No drama in it. Both resolved. The courtyard continues in its afternoon condition.
\end{action}

\begin{action}
\textsc{Rorario} does not look at the manuscript. He looks at the table. He looks at the window. He looks at nothing in particular, which is to say his gaze has ceased to be directed by any object of attention and is resting in the unfocused mode that precedes or follows sustained concentration, the visual system present but not engaged in any specific task.

The manuscript on the table is the minimal residue that the constraint process that constituted his diplomatic career has produced. It is not an expression of interior genius. It is not a conclusion arrived at. It is what remained after the accumulated eliminations of two decades --- what could still be said after everything that could not be said had been removed from the available option-space by the pressures of the institutional, the geopolitical, the perceptual, and the epistemic fields through which he has moved.

The pen is at the table's edge. The dogs are settled in the courtyard. The light will continue to decline. The manuscript will remain. End of State~III.
\end{action}

\transition{CUT TO:}

\transition{CUT TO:}

% ------------------------------------------------------------
\scenenumber{16c — The Ants}
% ------------------------------------------------------------

\sceneheading{EXT. Stone Wall, Garden — Day}

\begin{action}
The surface of a garden wall, close. The wall is old enough that its mortar has weathered into a texture of small channels and depressions, its face a topography at the scale of the creatures that traverse it daily --- a landscape whose features are real at that scale and invisible at the scale of the human figure passing beyond the frame of this observation. The scene occupies itself entirely with this surface and its traffic.

A column of ants moves along one of the wall's established routes, a path whose existence is the accumulated product of prior transit rather than any deliberate planning, the path itself being simply the set of points that prior transit has made chemically attractive to subsequent transit, the route self-reinforcing through the mechanism of its own use. The column moves with a consistency that is not mechanical but biological: individual ants deviate, investigate minor features of the surface, return to the column or rejoin it further along; the column as a whole maintains its trajectory and its approximate density through these individual variations in the way that a river maintains its course through the individual irregularity of each water molecule's path.
\end{action}

\begin{action}
At one point on the wall, a piece of food --- a fragment of something, its origin indeterminate from the available visual evidence, its nutritional relevance to the colony established by the behavior it elicits --- has been located by advance members of the column. The process of its transport begins. Individual ants arrive, assess, attach, lift, and begin to move. The fragment is substantially larger than any individual ant. Its transport is a collective problem whose solution requires the coordination of many individually insufficient forces into a combined effort sufficient to move the fragment along the route.

The coordination is not directed. No ant occupies a supervisory position from which it oversees and corrects the others. Each ant responds to the immediate mechanical situation it encounters --- the current position and orientation of the fragment, the forces being applied by adjacent ants, the surface features at the immediate point of contact --- and adjusts its own contribution accordingly. The collective movement that results from these local adjustments is neither the intention of any individual ant nor the product of any planning process, yet it is effective: the fragment moves along the route, the column reorganizes itself around the transport task, the wall's established path is being used for a purpose that its establishment did not anticipate and does not impede.

The scene offers nothing beyond this process, which continues for as long as the scene continues. No human figure enters the frame. The wall does not end within the available observation. The column continues.
\end{action}

\transition{CUT TO:}

% ------------------------------------------------------------
\scenenumber{16d — Night Insects}
% ------------------------------------------------------------

\sceneheading{EXT. Exterior Wall, Garden — Night}

\begin{action}
A lantern or torch has been fixed at a point on a wall, its light sufficient to illuminate a radius of perhaps two meters around the fixation point, beyond which the darkness is near-complete. The quality of this illumination is warm and unsteady, the flame's minor fluctuations producing continuous small changes in the distribution of light across the wall's surface and the air immediately in front of it.

Moths have arrived in the vicinity of the light source, drawn by mechanisms whose precise nature has been subject to various accounts across the period the film occupies, none of which constitutes the explanation that subsequent centuries would advance. They are present, which is the fact of the matter from inside the scene's available observation, and their behavior in proximity to the light exhibits the characteristic pattern: an approach along a spiral trajectory that brings them progressively closer to the source, the spiral's pitch varying between individuals and across time for the same individual, interrupted by brief departures and resumed approaches, the interruptions not terminating the overall movement toward proximity but only delaying it.
\end{action}

\begin{action}
The moths do not contact the flame during the scene's available duration. They are in continuous proximity to it, their trajectories passing near it repeatedly, the approach always arrested before completion and the withdrawal always followed by renewed approach. The behavior has the formal structure of an iteration that does not converge within any finite time --- each approach is closer to completion than the last in some respects and no closer in others, the system maintaining a condition of perpetual near-contact without achieving the contact toward which it appears to tend.

This is not failure in any evaluable sense. The moths are doing what moths do in proximity to light, which is to approach it and orbit it and approach it again, and the system constituted by their collective presence around the light source is stable: it will persist as long as the light persists and the moths persist, neither requiring resolution nor achieving it. A human figure moves through the lantern's light briefly, casting a shadow across the moths' flight space, then passes on. The moths redistribute slightly around the new configuration of light and darkness and resume their trajectories. The lantern continues to burn.
\end{action}

\transition{CUT TO:}

% ------------------------------------------------------------
\scenenumber{13a — Cain's Sign}
% ------------------------------------------------------------

\sceneheading{EXT. City Gate — Day}

\begin{action}
The gate is one of several points of controlled entry into a walled town, its architecture functional rather than ceremonial --- a passage through a wall, with the addition of the minimum apparatus required to manage who passes through it. At this hour, traffic is moderate; people and goods move in both directions through the passage, each entry and exit receiving the attention of the gate's attendants in proportion to its apparent relevance to their supervisory function, which is to say that most transit receives minimal attention and specific transit receives specific attention, the criteria distinguishing the two being partly visible from outside the attendants' reasoning and partly not.

A figure approaches the gate from the exterior road. He is of unremarkable aspect in most respects: his dress, his manner of movement, the goods or lack of goods he carries. He has a mark on his face --- on his forehead or cheek, its precise location established by the scene and not described here because the description would fix a specificity that the scene should not fix. The mark is neither fresh nor ancient; it has the character of something that has been present for long enough to be part of the figure's normal appearance rather than an event upon it.
\end{action}

\begin{action}
The gate attendants observe his approach. Their response to his arrival is, in one respect, continuous with their response to all arrivals: they register his presence and assess it. In another respect it differs: the presence of the mark produces in their response a quality that is not present in their response to unmarked arrivals, a quality that is difficult to characterize as either deference or avoidance because it contains elements of both without being reducible to either. They step back slightly. They do not challenge him. He passes through the gate without exchange.

A second attendant, who has been observing from a position further from the gate's passage, speaks briefly to the first after the marked figure has passed. The first responds. The exchange is too brief and too quiet to constitute anything that can be reconstructed from within the scene as a specific statement or a specific response. It registers only as an exchange --- two parties sharing something in the aftermath of an event, the content of the sharing not available for external determination. The marked figure continues along the road beyond the gate without looking back. He neither acknowledges the deference nor appears to be unaware of it. He proceeds.
\end{action}

\transition{CUT TO:}

% ------------------------------------------------------------
\scenenumber{9b — Abel's Altar Accident}
% ------------------------------------------------------------

\sceneheading{EXT. Open Ground, Hillside — Day}

\begin{action}
A stone surface of the kind used as an altar or offering place, its elevation above the ground modest and its construction informal --- stones gathered and placed rather than cut and fitted. Upon the surface, an arrangement of dry plant material and other organic matter has been set with some care, the arrangement suggesting a prior purpose whose execution is either complete or interrupted at this moment, the distinction between these two conditions not immediately resolvable from the scene's available evidence.

Two figures are visible at a distance from the stone surface, turned partially away from it, engaged in some task or conversation that has drawn their attention away from the surface itself. A third figure is nearer the surface, engaged with it in some capacity --- adjusting something, or examining it, or simply present in its vicinity. The wind moves through the hillside at intervals.
\end{action}

\begin{action}
At one such interval, the wind's passage across the stone surface displaces a portion of the arrangement placed upon it. The displacement is not violent; it has the character of a gradual shift rather than a sudden disturbance, several elements moving from their placed positions and coming to rest in new positions that the surface's slight unevenness and the wind's direction together determine. The third figure, near the surface, reaches to correct the displacement and in doing so introduces a further perturbation whose precise nature and consequence are not fully visible from the distance at which the observation is being conducted.

The two distant figures turn. They observe the surface in its current condition. The third figure is still beside it, his hand withdrawn, the arrangement now in a state that differs from its original arrangement in ways that could be attributed to several causes --- the wind, the corrective gesture, the perturbation introduced by the correction, or some combination of these acting in a sequence that is no longer reconstructable from the surface's current appearance. The three figures regard the surface. None speaks, within the available observation. The surface does not clarify what has occurred upon it. The wind continues at its intervals, moving through the hillside with the indifference of wind to the significances that surfaces acquire through their designation.
\end{action}

\transition{CUT TO:}

% ------------------------------------------------------------
\scenenumber{10b — The Staged Argument}
% ------------------------------------------------------------

\sceneheading{INT. or EXT. Performance Space — Evening}

\begin{action}
A space that has been prepared for a performance of the judicial or philosophical entertainment type: two positions marked by chairs or lecterns placed opposite one another, with a third position between and slightly elevated from them, occupied by a figure whose role is to receive, assess, and eventually determine the outcome of what the two opposing figures present. An audience occupies the surrounding space in the arrangement of bodies that performances of this kind generate --- proximate, attentive, with the particular quality of attention that entertainment whose outcome is ostensibly undecided produces, which is different from the attention produced by entertainment whose outcome is known in advance.

The entertainment has been underway for some time before the scene's entry. Each of the two advocates has had an opportunity to present their position, and both positions have been presented in sufficient detail that the audience has been given the materials to form preliminary assessments of relative merit, though the entertainment's form does not invite them to declare these assessments. The first advocate's position is internally coherent. The second advocate's position is internally coherent. They are not reconcilable.
\end{action}

\begin{action}
A third figure enters --- not the arbiter, who has been present throughout, but a figure entering from outside the performance's established cast, whose entry has the character of an interruption that the performance's form accommodates without having designed for it. He advances a third position. The third position is also internally coherent. It incorporates elements from each of the prior two positions while being reducible to neither, its relationship to its sources being one of selective extraction rather than synthesis --- the same operation that \textsc{Rorario} has been performing throughout the film in the administrative register, here externalized as a theatrical act. The audience receives this third position with the quality of attention that unanticipated developments in an entertainment produce: heightened, slightly uncertain, awaiting the structure's response to the interruption.

The arbiter deliberates. His deliberation is performed --- it has the visible marks of deliberation, the pauses and the expressions of consideration, without being demonstrably the actual process of the deliberation rather than its representation. He speaks. The verdict he delivers does not correspond to either of the two primary positions or to the third position; it is a formulation that acknowledges each while committing to none, a statement whose relationship to the material presented is one of formal reception rather than substantive adjudication. The audience's response is mixed in the specific manner that verdicts of this kind produce: those who had invested in one position experience the response appropriate to neither victory nor defeat, a condition for which the entertainment does not have a designated response, so each member of the audience resolves the condition individually according to their own resources. Some laugh. Some do not. \textsc{Rorario} is present as a spectator. He does not react to the verdict.
\end{action}

\transition{CUT TO:}

% ------------------------------------------------------------
\scenenumber{10c — The Puppet Theater}
% ------------------------------------------------------------

\sceneheading{EXT. Market or Public Space — Day}

\begin{action}
A traveling puppet show has established itself in a public space at the margin of a market, its booth of the standard construction: a frame that separates the operator's space from the performance space, with the performance space visible to the audience and the operator's space not, the convention of the form requiring the audience to maintain the pretense that the figures animating the puppets are the puppets themselves rather than the hands of an operator concealed immediately behind the visible plane.

The narrative being performed is of the heroic or scriptural type, its characters legible through their costuming and their relative scale as a story about a figure of authority, a figure of challenge, and a figure who intervenes between them. The \textsc{Narrator} --- positioned outside the booth, visible to the audience in a way that the operator is not --- speaks the story aloud as the puppets perform it. His narration and the puppets' actions are largely but not exactly coordinated.
\end{action}

\begin{action}
At a point in the performance where the Narrator describes an action --- an approach, a confrontation, a moment of decision --- the puppet representing the relevant figure performs a different action: not a contradictory action, but one that is adjacent to the narrated action in the story's logical space without being identical to it. The Narrator's account of what is happening and the puppet's enactment of what is happening are both coherent individually and slightly divergent from each other, the divergence of the kind that would be legible as an error if error were the available interpretation but that the performance's overall continuity makes available instead as a variation, a version, a rendering that is not wrong so much as differently decided.

The audience watches both the narration and the action, which is to say it watches a performance that is offering two slightly different accounts of the same event simultaneously, without marking the difference between them as requiring resolution. Some members of the audience follow the narration more closely; others follow the puppets. The performance does not indicate which is authoritative. It does not indicate that there is a distinction to be made. It simply continues, the narration and the action proceeding in their approximate parallel, each coherent, neither fully congruent with the other, the gap between them the space in which the audience's own determination of what the story is occurs, varying from one observer to the next without any of them being aware that they are not watching the same performance.

The performance ends. The Narrator concludes his account. The booth is closed. The audience disperses. The story that each member carries away is the story that their particular distribution of attention between the narration and the action has produced for them, which is not the same story for all of them, though none of them knows this.
\end{action}

\transition{CUT TO:}

% ============================================================
\statebreak{IV}{The Bayle Coda}
% ============================================================

\noindent{\itshape Spherepop operator: Quotient morphism. Late seventeenth century. The perceptual instability of the preceding states disappears entirely --- no flickers, no extended holds, no reweighted signal. Everything is correctly framed, correctly cut, properly legible. This is not a relief. It is the loss made visible through its own absence.}

\vspace{1.5em}

% ------------------------------------------------------------
\scenenumber{19 — The Printing House}
% ------------------------------------------------------------

\sceneheading{INT. Printing House, Amsterdam — Day}

\begin{action}
The printing house presents itself as a space of achieved order --- not the approximate order of the administrative chambers, whose surfaces bore the continuous evidence of processes incompletely concluded, but a more thoroughgoing order, one in which the relationship between function and arrangement has been stabilized through the repeated discipline of a craft that depends upon precision for its outputs. The type cases are organized. The press is calibrated. The surfaces are clean in the specific sense of surfaces from which what does not belong has been removed rather than surfaces that have never been used. It is the cleanliness of a system that knows what it is for.

The light enters from windows whose placement was determined by the requirements of fine work rather than by any ambient consideration, and it falls upon the working surfaces with the even, adequate illumination of light that has been put to use. Nothing here is underlit. Nothing is dramatically overlit. The room's illumination is, in the most functional sense possible, correct.
\end{action}

\begin{action}
A \textsc{Woman} who sets type moves through the print shop with the navigational fluency of someone for whom the space is so thoroughly known that its furniture and equipment register not as obstacles to be avoided but as features of the environment she inhabits as naturally as she inhabits her own body. She carries a small bundle of manuscript pages, their physical character immediately distinguishable from the standard copy that passes through the shop in the ordinary course of its operations: the paper is older, its surface bearing the particular texture of pages that have been handled many times across many years, its margins bearing faint additions and marginal notes in a hand different from the primary script. The pages have the physical biography of a document that was produced under circumstances --- various rooms, various inks, various degrees of pressure and haste --- that the document itself cannot narrate.

She places the bundle on the primary compositor's table. She reads one passage --- the gaze traveling across a few lines at functional reading speed rather than the contemplative speed of someone encountering unfamiliar material with sustained attention. She sets the bundle down. She continues to her next task. The interval between taking up the pages and setting them down again is perhaps twelve seconds. Nothing in her expression or subsequent behavior indicates that the interval has registered as different from the twelve seconds that preceded it or the twelve seconds that will follow it. She processes and continues. This is what processing without retention looks like from the outside.
\end{action}

\transition{CUT TO:}

% ------------------------------------------------------------
\scenenumber{20 — The Compositor}
% ------------------------------------------------------------

\sceneheading{INT. Printing House, Amsterdam — Day (continuous)}

\begin{action}
A \textsc{Compositor} receives the manuscript pages at his station. He is a man for whom text has become, through long professional acquaintance, a purely physical phenomenon --- a sequence of marks requiring transcription into the ordered material of type rather than a carrier of propositional content requiring evaluation. He reads the manuscript the way a skilled musician sight-reads a score: attending to the marks and their structural relations without necessarily inhabiting their semantic content at any depth. This is not incuriosity; it is the developed capacity of someone whose skill consists in the faithful reproduction of a text rather than its interpretation.

He works with the practiced efficiency of a man who has set many thousands of pages of type and for whom the individual decisions --- the selection of each sort from its case, the spacing of words, the management of the composing stick --- are executed below the threshold of deliberation, his hands finding the correct type while his eyes move ahead to the next clause. He does not pause to assess the quality of what he is setting.
\end{action}

\begin{action}
At one point he pauses at a word. Two variants are visible in the manuscript --- two different phrasings in different hands, one an apparent revision of the other. He selects one. He returns the type corresponding to the unselected variant to its compartment. The selection leaves no visible trace in the printed sequence; the discarded alternative is simply absent, not marked as absent.

The manuscript pages, set in type, become something structurally different from what they were. The pages had a physical character that corresponded to their history: different inks, different pressures, the uneven margins of pages produced across many years and many rooms. The type has no such history. Each sort is identical to every other sort of the same character; the spacing is regularized; the lines are justified to the measure of the column. Forty pages of manuscript become, through the compression of the entry format and the regularization of the type, eleven lines of set text. The compositor sets the eleventh line, places the final full stop, and slides the composed matter onto the galley. He takes up the next job. The manuscript pages remain at the edge of his station, their usefulness concluded.
\end{action}

\transition{CUT TO:}

% ------------------------------------------------------------
\scenenumber{21 — Bayle's Study}
% ------------------------------------------------------------

\sceneheading{INT. Study, Rotterdam — Day}

\begin{action}
Before the entry. \textsc{Pierre Bayle} at his working table, which is covered with the organized clutter appropriate to a man producing a reference work of encyclopedic scope: stacked volumes of varying ages and sizes, folded papers marking specific pages within them, a system of loose notes whose relation to the volumes they gloss is legible to Bayle and would require considerable reconstruction by anyone else. He works with the focused tranquility of someone engaged in a task that is intellectually demanding, procedurally clear, and adequately supplied with raw material.

He has before him three sources bearing information about Girolamo Rorario. Each source contains different information. Some of the information is consistent across two or more sources; some appears in only one. Some is inconsistent in ways that are irresolvable within the sources themselves --- different dates for the same event, different characterizations of the same position, references in one source to documents that the other sources do not mention and that have not been traced. Bayle reads each source with the careful attention of a scholar who knows that the material before him is imperfect and that the imperfection must be managed rather than concealed. He does not suppress the inconsistencies; he selects from among the available accounts those elements that are either consistent across sources or sufficiently singular in their attestation to merit inclusion as notable even without corroboration.
\end{action}

\begin{action}
The selection process is efficient. He reads a passage, determines its inclusion or exclusion in a matter of seconds, notes the former and passes by the latter with the practiced economy of a man who has made thousands of such decisions in the course of the \textit{Dictionnaire}'s production. What he excludes is not wrong, necessarily; it is unverifiable, or redundant, or insufficiently integrated with the entry's emerging structure to merit the space it would require. He sets aside letters whose dates do not align with the chronology he is constructing, references whose context he cannot reconstruct from the available material, characterizations of Rorario's positions that contradict the cleaner formulation he has found in another source and that would require extended footnoting to incorporate honestly.

The material he sets aside accumulates at the table's edge: a modest stack of pages and notes that represent what will not be in the entry. It is a larger stack than what will be in the entry. He does not regard it. He continues.

The entry grows toward its final form through a series of inclusions, each of which narrows the space available for the next. By the time it is nearly complete, its structure has become sufficiently determinate that the remaining decisions --- this phrase rather than that one, this date rather than the approximate one --- are constrained to a narrow range of admissible choices. The entry does not conclude so much as exhaust its admissible degrees of freedom. Bayle writes the final sentence. He reads it back once. He moves to the next entry.
\end{action}

\transition{CUT TO:}

% ------------------------------------------------------------
\scenenumber{22 — The Entry}
% ------------------------------------------------------------

\sceneheading{INT. Study, Rotterdam — Day}

\begin{action}
\textsc{Pierre Bayle} at his desk, the \textit{Dictionnaire historique et critique} in its production phase, the physical evidence of that production distributed across the surfaces of a room that has been organized for sustained intellectual labor: books accessible, writing materials replenished, light adequate, the ambient conditions managed to permit the kind of continuous concentrated work that the project requires. The room has the character of a space that has been submitted to a discipline --- not austere, but organized around a specific purpose to which its other potentials have been subordinated.

The entry for Rorario is complete. It occupies, in the manuscript before Bayle, eleven lines of text. It is accurate to the sources Bayle consulted. It is organized according to the principles of the \textit{Dictionnaire}'s format: the main text carrying the principal biographical and intellectual data, the footnotes providing the documentary apparatus and the qualifications that the main text's economy of space does not permit. It is legible. It is stable. It will remain substantially unchanged through the printed editions that follow this manuscript and will constitute, for the several centuries during which Rorario is known to history at all, the primary record of who he was and what he argued.

The provenance of the eleven lines is not marked within them. There is no notation indicating that the entry was produced from a particular set of sources, that certain sources were consulted and found inconsistent, that certain materials were set aside as unverifiable, that the manuscript from which the entry's intellectual content was derived had a physical history extending across two decades and several countries, that the man whose intellectual trajectory that manuscript recorded had watched a horse refuse a bridge, had held a pen above an incomplete sentence and set it down, had passed through a village where a dog pursued soap bubbles with total commitment. None of this is in the entry. Not because Bayle chose to exclude it --- it was never available to him as enterable content. It had not been transmitted through the channels that connect a life to its encyclopedic record. It had been the life, which is a different thing.
\end{action}

\begin{action}
The entry is legible. The entry is stable. Nothing in the entry indicates that anything is missing.

Bayle closes the manuscript volume and takes up the next. The \textit{Dictionnaire} continues its production. The entry for Rorario is complete within it, its eleven lines occupying the space that the entry format allocates, its text conforming to the conventions of the reference genre within which it exists, its information correct within the constraints of the sources from which it was derived. It is, within its own constraint set, a successful object.

The room is quiet. The light is adequate. The work continues. Outside, in the street below Bayle's window, the sounds of Rotterdam in the late seventeenth century constitute an environment that has no particular relation to the entry just completed and will continue to constitute that environment after the present scene concludes. The entry exists in the manuscript. It is stable. It will persist. It does not know what produced it. This is not a failure of the entry; it is the nature of entries.

The film ends here, on this: a room with adequate light, a scholar at his work, a completed entry in a manuscript that is otherwise continuing, and outside the window, a city that is going about its business in the manner of cities, which is to say without reference to any particular thing that has just been concluded within one of its buildings.

Nothing appears to have been lost.
\end{action}

% ============================================================
% END
% ============================================================

\end{document}
